% Chapter 1, Section 3 _Linear Algebra_ Jim Hefferon
%  http://joshua.smcvt.edu/linearalgebra
%  2001-Jun-09
\section{Bentuk Eselon Tereduksi}
Setelah mengembangkan mekanisme Metode Gauss,
kita melihat bahwa metode itu dapat dijalankan dengan lebih dari satu cara.
Sebagai contoh, dari matriks
\begin{equation*}
    \begin{mat}[r]
       2  &2  \\
       4  &3
    \end{mat}
\end{equation*}
kita dapat memperoleh salah satu dari ketiga matriks berbentuk eselon berikut.
\begin{equation*}
    \begin{mat}[r]
       2  &2  \\
       0  &-1
    \end{mat}
    \qquad
    \begin{mat}[r]
       1  &1  \\
       0  &-1
    \end{mat}
    \qquad
    \begin{mat}[r]
       2  &0  \\
       0  &-1
    \end{mat}
\end{equation*}
Matriks pertama dihasilkan oleh $-2\rho_1+\rho_2$.
Matriks kedua diperoleh dengan menjalankan $(1/2)\rho_1$, lalu
$-4\rho_1+\rho_2$.
Matriks ketiga diperoleh dari $-2\rho_1+\rho_2$, dilanjutkan dengan
$2\rho_2+\rho_1$ (setelah kombinasi baris pertama, matriks sudah berbentuk
eselon, tetapi operasi berikutnya tetap merupakan operasi baris yang sah).

Dalam bagian pertama bab ini, kita mencatat bahwa hal tersebut menimbulkan
beberapa pertanyaan.
Apakah sebarang dua versi berbentuk eselon dari suatu sistem linear memiliki
jumlah variabel bebas yang sama?
Jika ya, apakah keduanya memiliki variabel bebas yang persis sama?
Dalam bagian ini, kita akan memberikan cara untuk menentukan apakah suatu
sistem linear dapat diperoleh dari sistem lain melalui operasi baris.
Dari sana akan diperoleh jawaban ``ya'' untuk kedua pertanyaan tersebut.








\subsection{Reduksi Gauss--Jordan}%
% Gaussian elimination coupled with back-substitution
% solves linear systems but it is not the only method possible.
Berikut adalah perluasan Metode Gauss yang memiliki beberapa keunggulan.

\begin{example} \label{exm:GJRedReadOffSol}
Untuk menyelesaikan
\begin{equation*}
  \begin{linsys}{3}
    x  &+  &y  &-  &2z  &=  &-2  \\
       &   &y  &+  &3z  &=  &7   \\
    x  &   &   &-  &z   &=  &-1  
  \end{linsys}
\end{equation*}
kita dapat mulai seperti biasa dengan mereduksinya ke bentuk eselon.
\begin{equation*}
  \grstep{-\rho_1+\rho_3}
    \begin{amat}[r]{3}
       1  &1  &-2 &-2  \\
       0  &1  &3  &7   \\
       0  &-1 &1  &1
    \end{amat}
  \grstep{\rho_2+\rho_3}
    \begin{amat}[r]{3}
       1  &1  &-2 &-2  \\
       0  &1  &3  &7   \\
       0  &0  &4  &8
    \end{amat}
\end{equation*}
Kita dapat melanjutkan ke tahap kedua dengan menjadikan entri-entri utama
bernilai \( 1 \)
\begin{equation*}
    \grstep{(1/4)\rho_3}
    \begin{amat}[r]{3}
       1  &1  &-2 &-2  \\
       0  &1  &3  &7   \\
       0  &0  &1  &2
    \end{amat}
\end{equation*}
lalu ke tahap ketiga, yang menggunakan entri-entri utama untuk mengeliminasi
semua entri lain dalam setiap kolom melalui kombinasi ke arah atas.
\begin{equation*}
  \grstep[2\rho_3+\rho_1]{-3\rho_3+\rho_2}
    \begin{amat}[r]{3}
       1  &1  &0  &2   \\
       0  &1  &0  &1   \\
       0  &0  &1  &2
    \end{amat}
  \grstep{-\rho_2+\rho_1}
    \begin{amat}[r]{3}
       1  &0  &0  &1   \\
       0  &1  &0  &1   \\
       0  &0  &1  &2
    \end{amat}
\end{equation*}
Jawabannya adalah \( x=1 \), \( y=1 \), dan \( z=2 \).
\end{example}
%<*Pivoting>
Menggunakan satu entri untuk menolkan semua entri lain dalam suatu kolom
disebut \definend{pemivotan}\index{pemivotan} pada entri tersebut.
%</Pivoting>

Perhatikan bahwa operasi kombinasi baris pada tahap pertama bergerak dari kiri
ke kanan, sedangkan operasi kombinasi pada tahap ketiga bergerak dari kanan ke
kiri.

\begin{example}  \label{exm:GJRedReadOffSolTwo}
Operasi-operasi tahap tengah yang menjadikan entri-entri utama bernilai
\( 1 \) tidak saling berinteraksi, sehingga beberapa operasi dapat digabungkan
menjadi satu langkah.
\begin{align*}
    \begin{amat}[r]{2}
       2   &1   &7   \\
       4   &-2  &6
    \end{amat}
  &\grstep{-2\rho_1+\rho_2}
  \begin{amat}[r]{2}
       2   &1   &7   \\
       0   &-4  &-8
    \end{amat}                                   \\
  &\grstep[(-1/4)\rho_2]{(1/2)\rho_1}
  \begin{amat}[r]{2}
       1   &1/2   &7/2   \\
       0   &1     &2
    \end{amat}                                    \\
  &\grstep{-(1/2)\rho_2+\rho_1}
  \begin{amat}[r]{2}
       1   &0   &5/2   \\
       0   &1   &2
    \end{amat}
\end{align*}
Jawabannya adalah $x=5/2$ dan $y=2$.
\end{example}

%<*GaussJordanReduction>
Perluasan Metode Gauss ini disebut
\definend{Metode Gauss--Jordan}\index{Metode Gauss!Metode Gauss--Jordan} atau
\definend{reduksi Gauss--Jordan}.\index{persamaan linear!solusi!Gauss--Jordan}\index{Metode Gauss!Gauss--Jordan}
%</GaussJordanReduction>
% It goes past echelon form to a more refined, more specialized,
% matrix form.

\begin{definition}\label{def:RedEchForm}
%<*df:RedEchForm>
Suatu matriks atau sistem linear berada dalam
\definend{bentuk eselon tereduksi\/}\index{bentuk eselon!tereduksi}\index{bentuk eselon tereduksi}
jika, selain berada dalam bentuk eselon, setiap entri utama bernilai~$1$ dan
merupakan satu-satunya entri tak nol dalam kolomnya.
%</df:RedEchForm>
\end{definition}

\noindent
%<*CostRedEchForm>
Kerugian menggunakan reduksi Gauss--Jordan untuk menyelesaikan suatu sistem
adalah bertambahnya perhitungan aritmetika.
Keuntungannya adalah kita dapat langsung membaca deskripsi himpunan solusinya.
%</CostRedEchForm>

Pada setiap sistem berbentuk eselon, tereduksi ataupun tidak, kita dapat
langsung mengenali kapan sistem memiliki himpunan solusi kosong karena ada
persamaan yang kontradiktif.
Kita dapat mengenali kapan sistem memiliki himpunan solusi beranggota satu
karena tidak ada kontradiksi dan setiap variabel merupakan variabel utama pada
suatu baris.
Kita juga dapat mengenali kapan sistem memiliki himpunan solusi tak hingga
karena tidak ada kontradiksi dan sedikitnya satu variabel bebas.

Namun, dalam bentuk eselon tereduksi, kita tidak hanya dapat membaca ukuran
himpunan solusi, tetapi juga deskripsinya.
Tentu saja, tidak sulit mendeskripsikan himpunan solusi yang kosong.
\nearbyexample{exm:GJRedReadOffSol} dan~\ref{exm:GJRedReadOffSolTwo}
menunjukkan bahwa dalam kasus himpunan solusi beranggota satu, anggota tunggal
itu terdapat pada kolom konstanta.
Contoh berikut menunjukkan cara membaca parametrisasi himpunan solusi tak
hingga.

\begin{example}
\begin{multline*}
  \begin{amat}[r]{4}
     2  &6  &1  &2  &5  \\
     0  &3  &1  &4  &1  \\
     0  &3  &1  &2  &5
  \end{amat}
  \grstep{-\rho_2+\rho_3}
  \begin{amat}[r]{4}
     2  &6  &1  &2  &5  \\
     0  &3  &1  &4  &1  \\
     0  &0  &0  &-2 &4
  \end{amat}                                        \\
  \grstep[(1/3)\rho_2 \\ -(1/2)\rho_3]{(1/2)\rho_1}
  \repeatedgrstep[-\rho_3+\rho_1]{-(4/3)\rho_3+\rho_2}
  \repeatedgrstep{-3\rho_2+\rho_1}
  \begin{amat}[r]{4}
     1  &0  &-1/2  &0  &-9/2  \\
     0  &1  &1/3   &0  &3  \\
     0  &0  &0     &1  &-2
  \end{amat}
\end{multline*}
Sebagai sistem linear, matriks ini adalah
\begin{equation*}
  \begin{linsys}{4}
     x_1  &&     &-&1/2x_3  &&   &= &-9/2  \\
          &&x_2  &+&1/3x_3  &&   &= &3  \\
          &&     &&        &{}\hspace{.5em}{}&x_4 &= &-2    
  \end{linsys}
\end{equation*}
sehingga deskripsi himpunan solusinya adalah
\begin{equation*}  
  S=\set{\colvec{x_1 \\ x_2 \\ x_3 \\ x_4}
                        =\colvec[r]{-9/2 \\ 3 \\ 0 \\ -2}
                         +\colvec[r]{1/2 \\ -1/3 \\ 1 \\ 0}x_3
                        \suchthat x_3\in\Re}
\end{equation*}
\end{example}

Jadi, bentuk eselon bukanlah satu-satunya bentuk terbaik bagi semua sistem.
Bentuk lain, seperti bentuk eselon tereduksi, memiliki kelebihan dan
kekurangan.
Alih-alih membayangkan sistem linear (beserta matriks terkaitnya) sebagai objek
yang selalu kita operasikan menuju bentuk eselon, kita dapat memandang
sistem-sistem itu sebagai objek yang saling berhubungan dan dapat diubah dari
satu ke yang lain melalui operasi baris.
Sisa subbagian ini mengembangkan gagasan tersebut.

\begin{lemma} \label{le:RowOpsRev}
%<*lm:RowOpsRev>
Operasi baris elementer dapat dibalik.
%</lm:RowOpsRev>
\end{lemma}

\begin{proof}
%<*pf:RowOpsRev>
Untuk setiap matriks \( A \), efek menukar dua baris dibalik dengan menukarnya
kembali; mengalikan suatu baris dengan \( k \) tak nol dibatalkan dengan
mengalikannya dengan $1/k$; dan menambahkan kelipatan baris \( i \) ke baris
\( j \) (dengan $i\neq j$) dibatalkan dengan mengurangkan kelipatan baris
\( i \) yang sama dari baris \( j \).
\begin{equation*}
      A
     \grstep{\rho_i\leftrightarrow\rho_j}
     \repeatedgrstep{\rho_j\leftrightarrow\rho_i}
      A
  \qquad
        A
     \grstep{k\rho_i}
     \repeatedgrstep{(1/k)\rho_i}
      A
  \qquad
        A
     \grstep{k\rho_i+\rho_j}
     \repeatedgrstep{-k\rho_i+\rho_j}
      A                          
\end{equation*}
%</pf:RowOpsRev>
(Kita memerlukan syarat $i\neq j$; lihat
\nearbyexercise{exer:INotJMakesRowOpsRev}.)
\end{proof}

Sekali lagi, sudut pandang yang sedang kita kembangkan, dan kini didukung oleh
lema tersebut, adalah bahwa istilah `mereduksi menjadi' dapat menyesatkan:~jika
\( A\longrightarrow B \), kita tidak seharusnya menganggap \( B \) muncul
sesudah~\( A \) atau lebih sederhana daripada~$A$.
Sebaliknya, kita perlu memandang kedua matriks itu sebagai objek yang saling
berhubungan.
Gambar berikut menunjukkan matriks-matriks dari awal bagian ini beserta versi
bentuk eselon tereduksinya sebagai satu kelompok matriks yang dapat saling
direduksi.
\begin{center}  
  \includegraphics{gr/mp/ch1.28}
\end{center}

%<*EquivMatrices>
Kita mengatakan bahwa matriks-matriks yang dapat saling direduksi ekuivalen
terhadap relasi reduksi baris.
Hasil berikut membenarkan istilah tersebut dengan menggunakan definisi
ekuivalensi.\appendrefs{relasi ekuivalensi}
%</EquivMatrices>

\begin{lemma} \label{lm:ReducesToIsEqRel}
%<*lm:ReducesToIsEqRel>
Pada himpunan matriks, `mereduksi menjadi' merupakan relasi ekuivalensi.
%</lm:ReducesToIsEqRel>
\end{lemma}

\begin{proof}
%<*pf:ReducesToIsEqRel0>
Kita harus memeriksa syarat-syarat berikut:
(i)~refleksivitas, yaitu setiap matriks mereduksi menjadi dirinya sendiri;
(ii)~simetri, yaitu jika \( A \) mereduksi menjadi \( B \), maka
   \( B \) mereduksi menjadi \( A \); dan
(iii)~transitivitas, yaitu jika \( A \) mereduksi menjadi \( B \) dan
      \( B \) mereduksi menjadi \( C \), maka \( A \) mereduksi menjadi
      \( C \).
%</pf:ReducesToIsEqRel0>

%<*pf:ReducesToIsEqRel1>
Refleksivitas mudah; setiap matriks mereduksi menjadi dirinya sendiri melalui
nol operasi.

Relasi tersebut simetris menurut lema sebelumnya\Dash jika \( A \) mereduksi
menjadi \( B \) melalui sejumlah operasi baris, maka \( B \) juga mereduksi
menjadi \( A \) dengan membalik operasi-operasi itu.
%</pf:ReducesToIsEqRel1>

%<*pf:ReducesToIsEqRel2>
Untuk transitivitas, misalkan \( A \) mereduksi menjadi \( B \) dan
\( B \) mereduksi menjadi \( C \).
Meneruskan langkah-langkah reduksi $A \rightarrow\cdots\rightarrow B$
dengan langkah-langkah $B \rightarrow\cdots\rightarrow C$ menghasilkan
reduksi dari \( A \) ke \( C \).
%</pf:ReducesToIsEqRel2>
\end{proof}

\begin{definition} \label{df:RowEquivalence}
%<*df:RowEquivalence>
Dua matriks yang dapat saling direduksi melalui operasi baris elementer disebut
\definend{ekuivalen baris}.\index{matriks!ekuivalensi baris}%
\index{ekuivalensi baris}\index{relasi ekuivalensi!ekuivalensi baris}
%</df:RowEquivalence>
\end{definition}

%<*RowEquivalanceClasses>
Diagram berikut menunjukkan kumpulan semua matriks sebagai sebuah kotak.
Di dalam kotak itu, setiap matriks terletak dalam suatu kelas.
Dua matriks berada dalam kelas yang sama jika dan hanya jika keduanya dapat
saling direduksi.
Kelas-kelas tersebut saling lepas\Dash tidak ada matriks yang berada dalam dua
kelas berbeda.
Dengan demikian, kita telah mempartisi kumpulan matriks menjadi
\definend{kelas ekuivalensi baris}.\appendrefs{partisi dan wakil kelas}\index{partisi!kelas ekuivalensi baris}
%</RowEquivalanceClasses>
\begin{center}
  \includegraphics{gr/mp/ch1.27}
\end{center}
\noindent Salah satu kelas adalah kelompok matriks yang saling berhubungan
dari awal bagian ini, sebagaimana digambarkan di atas (kelas tersebut memuat
semua matriks nonsingular $\nbyn{2}$).

Subbagian berikut membuktikan bahwa bentuk eselon tereduksi suatu matriks
bersifat unik.
Dalam bahasa relasi ekuivalensi baris, kita akan membuktikan bahwa setiap
matriks ekuivalen baris dengan tepat satu matriks berbentuk eselon tereduksi.
Dalam bahasa partisi, kita akan membuktikan bahwa setiap kelas ekuivalensi
memuat tepat satu matriks berbentuk eselon tereduksi.
Jadi, setiap matriks berbentuk eselon tereduksi menjadi wakil kelasnya.

\begin{exercises}
   \recommended \item 
     Gunakan reduksi Gauss--Jordan untuk menyelesaikan setiap sistem.
     \begin{exparts*}
        \partsitem \(
          \begin{linsys}[t]{2}
               x  &+  &y  &=  &2  \\
               x  &-  &y  &=  &0  
          \end{linsys}   \)
        \partsitem \(
          \begin{linsys}[t]{3}
               x  &   &   &-  &z  &=  &4  \\
              2x  &+  &2y &   &   &=  &1  
          \end{linsys}  \)
        \partsitem  \(
           \begin{linsys}[t]{2}
               3x  &-  &2y  &=  &1  \\
               6x  &+  &y   &=  &1/2 
           \end{linsys}  \)
        \partsitem \(
           \begin{linsys}[t]{3}
              2x  &-  &y  &  &  &= &-1  \\
               x  &+  &3y &- &z &= &5   \\
                  &   &y  &+ &2z&= &5   
           \end{linsys} \)
     \end{exparts*}
     \begin{answer}
       Jawaban-jawaban ini hanya menunjukkan reduksi Gauss--Jordan.
       Setelah itu, deskripsi himpunan solusinya dapat dibaca dengan mudah.
       \begin{exparts}
         \partsitem Himpunan solusinya hanya memuat satu anggota.
           \begin{multline*}
             \begin{amat}[r]{2}
               1  &1  &2  \\
               1  &-1 &0
             \end{amat}
             \grstep{-\rho_1+\rho_2}
             \begin{amat}[r]{2}
               1  &1  &2  \\
               0  &-2 &-2
             \end{amat}                          \\                       
             \grstep{-(1/2)\rho_2}
             \begin{amat}[r]{2}
               1  &1  &2  \\
               0  &1  &1
             \end{amat}                         
             \grstep{-\rho_2+\rho_1}
             \begin{amat}[r]{2}
               1  &0  &1  \\
               0  &1  &1
             \end{amat}
           \end{multline*}
         \partsitem Himpunan solusinya memiliki satu parameter.
            \begin{equation*}
             \begin{amat}[r]{3}
               1  &0  &-1  &4  \\
               2  &2  &0   &1
             \end{amat}
             \grstep{-2\rho_1+\rho_2}
             \begin{amat}[r]{3}
               1  &0  &-1  &4  \\
               0  &2  &2   &-7
             \end{amat}                    
             \grstep{(1/2)\rho_2}
             \begin{amat}[r]{3}
               1  &0  &-1  &4  \\
               0  &1  &1   &-7/2
             \end{amat}
           \end{equation*}
         \partsitem Terdapat solusi tunggal.
           \begin{multline*}
             \begin{amat}[r]{2}
               3  &-2  &1  \\
               6  &1   &1/2
             \end{amat}
             \grstep{-2\rho_1+\rho_2}
             \begin{amat}[r]{2}
               3  &-2  &1  \\
               0  &5   &-3/2
             \end{amat}                            \\                       
             \grstep[(1/5)\rho_2]{(1/3)\rho_1}
             \begin{amat}[r]{2}
               1  &-2/3&1/3 \\
               0  &1   &-3/10
             \end{amat}                       
             \grstep{(2/3)\rho_2+\rho_1}
             \begin{amat}[r]{2}
               1  &0   &2/15 \\
               0  &1   &-3/10
             \end{amat}
          \end{multline*}
        \partsitem Pertukaran baris pada langkah kedua mempermudah
          perhitungan aritmetika.
         \begin{multline*}
          \begin{amat}[r]{3}
            2  &-1  &0  &-1  \\
            1  &3   &-1 &5   \\
            0  &1   &2  &5
          \end{amat}
          \grstep{-(1/2)\rho_1+\rho_2}
          \begin{amat}[r]{3}
            2  &-1  &0  &-1   \\
            0  &7/2 &-1 &11/2 \\
            0  &1   &2  &5
          \end{amat}                                \\
          \begin{aligned}
          &\grstep{\rho_2\leftrightarrow\rho_3}
          \begin{amat}[r]{3}
            2  &-1  &0  &-1   \\
            0  &1   &2  &5    \\
            0  &7/2 &-1 &11/2
          \end{amat}                      
            \grstep{-(7/2)\rho_2+\rho_3}
            \begin{amat}[r]{3}
              2  &-1  &0  &-1   \\
              0  &1   &2  &5    \\
              0  &0   &-8 &-12
            \end{amat}                               \\     
            &\grstep[-(1/8)\rho_3]{(1/2)\rho_1}
            \begin{amat}[r]{3}
              1  &-1/2&0  &-1/2 \\
              0  &1   &2  &5    \\
              0  &0   &1  &3/2
            \end{amat}                                
            \grstep{-2\rho_3+\rho_2}
            \begin{amat}[r]{3}
              1  &-1/2&0  &-1/2 \\
              0  &1   &0  &2    \\
              0  &0   &1  &3/2
            \end{amat}                                    \\             
            &\grstep{(1/2)\rho_2+\rho_1}
            \begin{amat}[r]{3}
              1  &0   &0  &1/2  \\
              0  &1   &0  &2    \\
              0  &0   &1  &3/2
            \end{amat}
          \end{aligned}
        \end{multline*}
       \end{exparts}  
     \end{answer}
  \item Lakukan reduksi Gauss--Jordan.
    \begin{exparts*}
      \partsitem
        $\begin{linsys}[t]{3}
          x  &+ &y &- &z &= &3 \\
          2x &- &y  &- &z &= &1 \\
          3x &+  &y  &+ &2z &= &0
        \end{linsys}$
      \partsitem
        $\begin{linsys}[t]{3}
          x  &+ &y  &+ &2z  &= &0 \\
          2x  &- &y  &+  &z &= &1 \\
          4x &+ &y  &+ &5z &= &1  
        \end{linsys}$
    \end{exparts*}
    \begin{answer}
      \begin{exparts}
        \partsitem
          \begin{multline*}
            \begin{amat}{3}
              1 &1  &-1 &3 \\
              2 &-1 &-1 &1 \\
              3 &1  &2  &0
            \end{amat}
            \grstep[-3\rho_1+\rho_3]{-2\rho_1+\rho_2}
            \begin{amat}{3}
              1 &1  &-1 &3 \\
              0 &-3 &1  &-5 \\
              0 &-2  &5  &-9
            \end{amat}                             \\
          \begin{aligned}   
            &\grstep{-(2/3)\rho_2+\rho_3}
            \begin{amat}{3}
              1 &1  &-1    &3 \\
              0 &-3 &1     &-5 \\
              0 &0  &13/3  &-17/3
            \end{amat}                                 
            \grstep[(3/13)\rho_3]{-(1/3)\rho_2}
            \begin{amat}{3}
              1 &1  &-1    &3 \\
              0 &1  &-1/3    &5/3 \\
              0 &0  &1      &-17/13
            \end{amat}                                \\
            &\grstep[(1/3)\rho_3+\rho_2]{\rho_3+\rho_1}
            \begin{amat}{3}
              1 &1  &0    &22/13 \\
              0 &1  &0    &16/13 \\
              0 &0  &1      &-17/13
            \end{amat}
            \grstep{-\rho_2+\rho_1}
            \begin{amat}{3}
              1 &0  &0    &6/13 \\
              0 &1  &0    &16/13 \\
              0 &0  &1      &-17/13
            \end{amat}
          \end{aligned}
          \end{multline*}
        \partsitem
          \begin{multline*}
            \begin{amat}{3}
              1 &1  &2 &0 \\
              2 &-1 &1 &1 \\
              4 &1  &5 &1
            \end{amat}
            \grstep[-4\rho_1+\rho_3]{-2\rho_1+\rho_2}
            \begin{amat}{3}
              1 &1  &2  &0 \\
              0 &-3 &-3 &1 \\
              0 &-3 &-3 &1
            \end{amat}
            \grstep{-\rho_2+\rho_3}
            \begin{amat}{3}
              1 &1  &2  &0 \\
              0 &-3 &-3 &1 \\
              0 &0  &0  &0
            \end{amat}                              \\
            \grstep{-(1/3)\rho_2}
            \begin{amat}{3}
              1 &1  &2  &0 \\
              0 &1  &1 &-1/3 \\
              0 &0  &0  &0
            \end{amat}
            \grstep{-\rho_2+\rho_1}
            \begin{amat}{3}
              1 &0  &1  &1/3 \\
              0 &1  &1 &-1/3 \\
              0 &0  &0  &0
            \end{amat}
          \end{multline*}
      \end{exparts}
    \end{answer}
  \recommended \item 
    Tentukan bentuk eselon tereduksi setiap matriks.
    \begin{exparts*}
      \partsitem \( \begin{mat}[r]
          2  &1  \\
          1  &3
        \end{mat}  \)
      \partsitem \( \begin{mat}[r]
          1  &3  &1  \\
          2  &0  &4  \\
         -1  &-3 &-3
        \end{mat}  \)
      \partsitem \( \begin{mat}[r]
          1  &0  &3  &1  &2  \\
          1  &4  &2  &1  &5  \\
          3  &4  &8  &1  &2
        \end{mat}  \)
      \partsitem \( \begin{mat}[r]
          0  &1  &3  &2  \\
          0  &0  &5  &6  \\
          1  &5  &1  &5
        \end{mat}  \)
    \end{exparts*}
    \begin{answer}
      Gunakan reduksi Gauss--Jordan.
      \begin{exparts}
        \partsitem Bentuk eselon tereduksinya hanya memuat nol, kecuali
          diagonal yang seluruh entrinya satu.
          \begin{equation*}
            \grstep{-(1/2)\rho_1+\rho_2}
            \begin{mat}[r]
              2  &1  \\
              0  &5/2
            \end{mat}
            \grstep[(2/5)\rho_2]{(1/2)\rho_1}
            \begin{mat}[r]
              1  &1/2\\
              0  &1
            \end{mat}                      
            \grstep{-(1/2)\rho_2+\rho_1}
            \begin{mat}[r]
              1  &0  \\
              0  &1
            \end{mat}
          \end{equation*}
        \partsitem Seperti pada soal sebelumnya, bentuk eselon tereduksinya
          hanya memuat nol, kecuali diagonal yang seluruh entrinya satu.
          \begin{multline*}
            \grstep[\rho_1+\rho_3]{-2\rho_1+\rho_2}
            \begin{mat}[r]
              1  &3  &1  \\
              0  &-6 &2  \\
              0  &0  &-2
            \end{mat}
            \grstep[-(1/2)\rho_3]{-(1/6)\rho_2}
            \begin{mat}[r]
              1  &3  &1     \\
              0  &1  &-1/3  \\
              0  &0  &1
            \end{mat}                                      \\              
            \grstep[-\rho_3+\rho_1]{(1/3)\rho_3+\rho_2}
            \begin{mat}[r]
              1  &3  &0     \\
              0  &1  &0     \\
              0  &0  &1
            \end{mat}                  
            \grstep{-3\rho_2+\rho_1}
            \begin{mat}[r]
              1  &0  &0     \\
              0  &1  &0     \\
              0  &0  &1
            \end{mat}
           \end{multline*}
        \partsitem Karena jumlah kolom lebih banyak daripada jumlah baris,
          hasilnya tidak mungkin hanya berupa diagonal satu.
          \begin{multline*}
            \grstep[-3\rho_1+\rho_3]{-\rho_1+\rho_2}
            \begin{mat}[r]
              1  &0  &3  &1  &2  \\
              0  &4  &-1 &0  &3  \\
              0  &4  &-1 &-2 &-4
            \end{mat}
            \grstep{-\rho_2+\rho_3}
            \begin{mat}[r]
              1  &0  &3  &1  &2  \\
              0  &4  &-1 &0  &3  \\
              0  &0  &0  &-2 &-7
            \end{mat}                           \\
            \grstep[-(1/2)\rho_3]{(1/4)\rho_2}
            \begin{mat}[r]
              1  &0  &3    &1  &2  \\
              0  &1  &-1/4 &0  &3/4  \\
              0  &0  &0    &1  &7/2
            \end{mat}                           
            \grstep{-\rho_3+\rho_1}
            \begin{mat}[r]
              1  &0  &3    &0  &-3/2  \\
              0  &1  &-1/4 &0  &3/4     \\
              0  &0  &0    &1  &7/2
            \end{mat}
          \end{multline*}
        \partsitem Seperti pada butir sebelumnya, matriks ini bukan matriks
          persegi.
          \begin{multline*}
            \grstep{\rho_1\leftrightarrow\rho_3}
            \begin{mat}[r]
              1  &5  &1  &5  \\
              0  &0  &5  &6  \\
              0  &1  &3  &2
            \end{mat}
            \grstep{\rho_2\leftrightarrow\rho_3}
            \begin{mat}[r]
              1  &5  &1  &5  \\
              0  &1  &3  &2  \\
              0  &0  &5  &6
            \end{mat}                    \\
            \begin{aligned}
              &\grstep{(1/5)\rho_3}
              \begin{mat}[r]
                1  &5  &1  &5  \\
                0  &1  &3  &2  \\
                0  &0  &1  &6/5
              \end{mat}                  
              \grstep[-\rho_3+\rho_1]{-3\rho_3+\rho_2}
              \begin{mat}[r]
                1  &5  &0  &19/5  \\
                0  &1  &0  &-8/5  \\
                0  &0  &1  &6/5
              \end{mat}                   \\
              &\grstep{-5\rho_2+\rho_1}
              \begin{mat}[r]
                1  &0  &0  &59/5  \\
                0  &1  &0  &-8/5  \\
                0  &0  &1  &6/5
              \end{mat}
            \end{aligned}
          \end{multline*}
      \end{exparts}  
    \end{answer}
  \item Tentukan bentuk eselon tereduksi masing-masing matriks berikut.
    \begin{exparts*}
      \partsitem $
        \begin{mat}
          0  &2 &1 \\
          2  &-1 &1 \\
          -2 &-1 &0
        \end{mat}$
      \partsitem $
        \begin{mat}
          1 &3 &1 \\
          2 &6 &2 \\
         -1 &0 &0
        \end{mat}$
    \end{exparts*}
    \begin{answer}
      \begin{exparts}
        \partsitem Mulailah dengan pertukaran baris.
          \begin{equation*}
            \grstep{\rho_1\leftrightarrow\rho_2}
            \grstep{\rho_1+\rho_3}
            \grstep{\rho_2+\rho_3}
            \grstep[(1/2)\rho_2 \\ (1/2)\rho_3]{(1/2)\rho_1}
            \grstep[(-1/2)\rho_3+\rho_2]{(-1/2)\rho_3+\rho_1}
            \grstep{(1/2)\rho_2+\rho_1}
            \begin{mat}
              1 &0 &0 \\
              0 &1 &0 \\
              0 &0 &1
            \end{mat}
          \end{equation*}
        \partsitem Di sini pertukaran baris dilakukan di tengah proses.
        \begin{equation*}
          \grstep[\rho_1+\rho_3]{-2\rho_1+\rho_2}
          \grstep{\rho_2\leftrightarrow\rho_3}
          \grstep{(1/3)\rho_2}
          \grstep{-3\rho_2+\rho_1}
          \begin{mat}
            1 &0 &0   \\
            0 &1 &1/3 \\
            0 &0 &0
          \end{mat}
        \end{equation*}
      \end{exparts}
    \end{answer}
  \recommended \item 
    Tentukan setiap himpunan solusi dengan menggunakan reduksi Gauss--Jordan,
    lalu baca parametrisasi yang dihasilkan.
    \begin{exparts}
      \partsitem \( \begin{linsys}[t]{3}
                  2x  &+  &y  &-  &z  &=  &1  \\
                  4x  &-  &y  &   &   &=  &3  
                  \end{linsys}  \)
      \partsitem \( \begin{linsys}[t]{4}
                   x  &   &   &-  &z  &   &   &=  &1  \\
                      &   &y  &+  &2z &-  &w  &=  &3  \\
                   x  &+  &2y &+  &3z &-  &w  &=  &7  
                    \end{linsys}  \)
      \partsitem \( \begin{linsys}[t]{4}
                   x  &-  &y  &+  &z  &   &   &=  &0  \\
                      &   &y  &   &   &+  &w  &=  &0  \\
                  3x  &-  &2y &+  &3z &+  &w  &=  &0  \\
                      &   &-y &   &   &-  &w  &=  &0  
                  \end{linsys}  \)
      \partsitem \( \begin{linsys}[t]{5}
                   a  &+  &2b &+  &3c &+  &d  &-  &e  &=  &1  \\
                  3a  &-  &b  &+  &c  &+  &d  &+  &e  &=  &3  
                  \end{linsys}  \)
    \end{exparts}
    \begin{answer}
      Untuk bagian reduksi Gauss, lihat jawaban Latihan~I.2 di Bab Satu,
      \nearbyexercise{exer:SlvMatNot}.
      \begin{exparts}
      \partsitem Bagian ``Jordan'' berjalan sebagai berikut.
        \begin{equation*}
          \grstep[-(1/3)\rho_2]{(1/2)\rho_1}
          \begin{amat}[r]{3}
            1  &1/2 &-1/2 &1/2  \\
            0  &1   &-2/3 &-1/3
          \end{amat}                           
          \grstep{-(1/2)\rho_2+\rho_1}
          \begin{amat}[r]{3}
            1  &0   &-1/6 &2/3  \\
            0  &1   &-2/3 &-1/3
          \end{amat}
        \end{equation*}
        Himpunan solusinya adalah
        \begin{equation*}
          \set{\colvec[r]{2/3 \\ -1/3 \\ 0}
               +\colvec[r]{1/6 \\ 2/3 \\ 1}z
              \suchthat z\in\Re}
        \end{equation*}
      \partsitem Bagian keduanya adalah
        \begin{equation*}
          \grstep{\rho_3+\rho_2}
          \begin{amat}[r]{4}
            1  &0  &-1  &0  &1 \\
            0  &1  &2   &0  &3 \\
            0  &0  &0   &1  &0
          \end{amat}
        \end{equation*}
        sehingga himpunan solusinya adalah
        \begin{equation*}
          \set{\colvec[r]{1 \\ 3 \\ 0 \\ 0}
               +\colvec[r]{1 \\ -2 \\ 1 \\ 0}z
              \suchthat z\in\Re}
        \end{equation*}
      \partsitem Bagian Jordan berikut
        \begin{equation*}
          \grstep{\rho_2+\rho_1}
          \begin{amat}[r]{4}
            1  &0  &1   &1  &0 \\
            0  &1  &0   &1  &0 \\
            0  &0  &0   &0  &0 \\
            0  &0  &0   &0  &0
          \end{amat}
        \end{equation*}
        menghasilkan
        \begin{equation*}
          \set{\colvec[r]{0 \\ 0 \\ 0 \\ 0}
               +\colvec[r]{-1 \\ 0 \\ 1 \\ 0}z
               +\colvec[r]{-1 \\ -1 \\ 0 \\ 1}w
              \suchthat z,w\in\Re}
        \end{equation*}
        (tentu saja, vektor nol dapat dihilangkan dari deskripsi tersebut).
      \partsitem Bagian ``Jordan''
        \begin{align*}
          &\grstep{-(1/7)\rho_2}
          \begin{amat}[r]{5}
            1  &2  &3   &1   &-1   &1  \\
            0  &1  &8/7 &2/7 &-4/7 &0
          \end{amat}                   \\
          &\grstep{-2\rho_2+\rho_1}
          \begin{amat}[r]{5}
            1  &0  &5/7 &3/7 &1/7  &1  \\
            0  &1  &8/7 &2/7 &-4/7 &0
          \end{amat}
        \end{align*}
        berakhir dengan himpunan solusi berikut.
        \begin{equation*}
          \set{\colvec[r]{1 \\ 0 \\ 0 \\ 0 \\ 0}
               +\colvec[r]{-5/7 \\ -8/7 \\ 1 \\ 0 \\ 0}c
               +\colvec[r]{-3/7 \\ -2/7 \\ 0 \\ 1 \\ 0}d
               +\colvec[r]{-1/7 \\ 4/7 \\ 0 \\ 0 \\ 1}e
              \suchthat c,d,e\in\Re}
        \end{equation*}
    \end{exparts}
   \end{answer}
  \item 
    Berikan dua versi berbentuk eselon yang berbeda dari matriks berikut.
    \begin{equation*}
      \begin{mat}[r]
        2  &1  &1  &3  \\
        6  &4  &1  &2  \\
        1  &5  &1  &5
      \end{mat}
    \end{equation*}
    \begin{answer}
      Metode Gauss yang biasa menghasilkan salah satunya:
      \begin{equation*}
        \grstep[-(1/2)\rho_1+\rho_3]{-3\rho_1+\rho_2}
        \begin{mat}[r]
          2  &1  &1  &3  \\
          0  &1  &-2 &-7 \\
          0  &9/2&1/2&7/2
        \end{mat}
        \grstep{-(9/2)\rho_2+\rho_3}
        \begin{mat}[r]
          2  &1  &1    &3  \\
          0  &1  &-2   &-7 \\
          0  &0  &19/2 &35
        \end{mat}
      \end{equation*}
      dan perubahan kosmetik apa pun, seperti mengalikan baris terbawah dengan
      \( 2 \),
      \begin{equation*}
        \begin{mat}[r]
          2  &1  &1    &3  \\
          0  &1  &-2   &-7 \\
          0  &0  &19   &70
        \end{mat}
      \end{equation*}
      menghasilkan versi lainnya.  
    \end{answer}
  \recommended \item \label{exer:PossRedEchFrms} 
    Daftarkan semua bentuk eselon tereduksi yang mungkin untuk setiap ukuran.
    \begin{exparts*}
      \partsitem \( \nbyn{2} \)
      \partsitem \( \nbym{2}{3} \)
      \partsitem \( \nbym{3}{2} \)
      \partsitem \( \nbyn{3} \)
    \end{exparts*}
    \begin{answer}
      Dalam kasus-kasus berikut, kita mengambil $a,b\in\Re$.
      Jadi, beberapa bentuk kanonis yang didaftarkan sebenarnya mencakup tak
      hingga banyak kasus.
      Secara khusus, bentuk-bentuk tersebut mencakup kasus $a=0$ dan $b=0$.
      \begin{exparts}
        \partsitem 
          $\begin{mat}[r]
            0  &0  \\
            0  &0
          \end{mat}$,
          $\begin{mat}[r]
            1  &a  \\
            0  &0
          \end{mat}$, 
          $\begin{mat}[r]
            0  &1  \\
            0  &0
          \end{mat}$, 
          $\begin{mat}[r]
            1  &0  \\
            0  &1
          \end{mat}$
        \partsitem
          $\begin{mat}[r]
               0  &0  &0  \\
               0  &0  &0
             \end{mat}$,
          $\begin{mat}[r]
               1  &a  &b  \\
               0  &0  &0
             \end{mat}$,
          $\begin{mat}[r]
               0  &1  &a  \\
               0  &0  &0
             \end{mat}$,
          $\begin{mat}[r]
               0  &0  &1  \\
               0  &0  &0
             \end{mat}$,
          $\begin{mat}[r]
               1  &0  &a  \\
               0  &1  &b
             \end{mat}$,
          $\begin{mat}[r]
               1  &a  &0  \\
               0  &0  &1
             \end{mat}$,
          $\begin{mat}[r]
               0  &1  &0  \\
               0  &0  &1
             \end{mat}$
        \partsitem
          $\begin{mat}[r]
               0  &0  \\
               0  &0  \\
               0  &0
             \end{mat}$,
          $\begin{mat}[r]
               1  &a  \\
               0  &0  \\
               0  &0
             \end{mat}$,
          $\begin{mat}[r]
               0  &1  \\
               0  &0  \\
               0  &0
             \end{mat}$,
          $\begin{mat}[r]
               1  &0  \\
               0  &1  \\
               0  &0
             \end{mat}$
        \partsitem
          $\begin{mat}[r]
               0  &0  &0  \\
               0  &0  &0  \\
               0  &0  &0
             \end{mat}$,
          $\begin{mat}[r]
               1  &a  &b  \\
               0  &0  &0  \\
               0  &0  &0
             \end{mat}$,
          $\begin{mat}[r]
               0  &1  &a  \\
               0  &0  &0  \\
               0  &0  &0
             \end{mat}$,
          $\begin{mat}[r]
               0  &1  &0  \\
               0  &0  &1  \\
               0  &0  &0
             \end{mat}$,
          $\begin{mat}[r]
               0  &0  &1  \\
               0  &0  &0  \\
               0  &0  &0
             \end{mat}$,
          $\begin{mat}[r]
               1  &0  &a  \\
               0  &1  &b  \\
               0  &0  &0
             \end{mat}$,
          $\begin{mat}[r]
               1  &a  &0  \\
               0  &0  &1  \\
               0  &0  &0
             \end{mat}$,
          $\begin{mat}[r]
               1  &0  &0  \\
               0  &1  &0  \\
               0  &0  &1
             \end{mat}$
      \end{exparts}  
    \end{answer}
  \recommended \item  
    Apa hasil penerapan reduksi Gauss--Jordan pada matriks nonsingular?
    \begin{answer}
      Sistem linear homogen nonsingular memiliki solusi tunggal.
      Jadi, matriks nonsingular harus tereduksi menjadi matriks (persegi) yang
      seluruh entrinya \( 0 \), kecuali entri-entri \( 1 \) pada diagonal dari
      kiri atas ke kanan bawah, seperti berikut.
      \begin{equation*}
         \begin{mat}[r]
           1  &0  \\
           0  &1  \\
         \end{mat}
         \qquad
         \begin{mat}[r]
           1  &0  &0  \\
           0  &1  &0  \\
           0  &0  &1
         \end{mat}
      \end{equation*}  
    \end{answer}
 \item Tentukan apakah setiap relasi berikut merupakan relasi ekuivalensi pada
   himpunan matriks $\nbyn{2}$.
   \begin{exparts}
     \partsitem dua matriks berelasi jika keduanya memiliki entri yang sama
       pada baris pertama dan kolom pertama
     % \partsitem two matrices are related if their four entries sum to the 
     %   same total
     \partsitem dua matriks berelasi jika keduanya memiliki entri yang sama
       pada baris pertama dan kolom pertama, atau entri yang sama pada baris
       kedua dan kolom kedua
   \end{exparts}
   \begin{answer}
     \begin{exparts}
       \partsitem Ini merupakan relasi ekuivalensi.
        
          Kita dapat menuliskan $M_1\sim M_2$ jika keduanya berelasi.
          Relasi $\sim$ refleksif karena setiap matriks memiliki entri $1,1$
          yang sama dengan dirinya sendiri.
          Relasi tersebut simetris karena jika $M_1$ memiliki entri $1,1$
          yang sama dengan~$M_2$, maka jelas $M_2$ juga memiliki entri $1,1$
          yang sama dengan~$M_1$.
          Terakhir, relasi tersebut transitif: jika $M_1\sim M_2$, keduanya
          memiliki entri $1,1$ yang sama; jika $M_2\sim M_3$, keduanya juga
          memiliki entri itu yang sama. Jadi, ketiganya memiliki entri $1,1$
          yang sama dan $M_1\sim M_3$.
       % \partsitem  This is an equivalence.
       %    Write $M_1\sim M_2$ if they are related.

       %    This relation is reflexive because any matrix has the 
       %    same sum of entries as itself.
       %    The $\sim$~relation is symmetric because if $M_1$ has the same
       %    sum of 
       %    entries as~$M_2$, then also $M_2$ has the same sum of entries
       %    as~$M_1$.
       %    Finally, the relation is transitive because if $M_1\sim M_2$ so their
       %    entry sum is the same, 
       %    and $M_2\sim M_3$ so they have the 
       %    same entry sum as each other, then all three have the same entry sum, 
       %    and $M_1\sim M_3$. 
       \partsitem
          Ini bukan relasi ekuivalensi karena tidak transitif.
          Matriks pertama dan kedua berikut berelasi melalui entri $1,1$,
          sedangkan matriks kedua dan ketiga berelasi melalui entri $2,2$.
          Namun, matriks pertama dan ketiga tidak berelasi.
          \begin{equation*}
            \begin{mat}
              1 &0 \\
              0 &0
            \end{mat}
            \quad
            \begin{mat}
              1 &0 \\
              0 &-1
            \end{mat}
            \quad
            \begin{mat}
              0 &0 \\
              0 &-1
            \end{mat}
          \end{equation*}
     \end{exparts}
   \end{answer}
 \item \cite{Cleary}
    Tinjau relasi berikut pada himpunan matriks $\nbyn{2}$: kita mengatakan
    bahwa $A$ \textit{serupa-jumlah} dengan $B$ jika jumlah semua entri dalam
    $A$ sama dengan jumlah semua entri dalam $B$.
    Sebagai contoh, matriks nol serupa-jumlah dengan matriks yang baris
    pertamanya memuat dua angka tujuh dan baris keduanya memuat dua angka
    negatif tujuh.
    Buktikan atau sangkal bahwa relasi ini merupakan relasi ekuivalensi pada
    himpunan matriks $\nbyn{2}$.
    \begin{answer}
      Relasi ini merupakan relasi ekuivalensi.
      Untuk membuktikannya, kita harus memeriksa bahwa relasi tersebut
      refleksif, simetris, dan transitif.

      Andaikan semua matriks berukuran $\nbyn{2}$.
      Untuk refleksivitas, perhatikan bahwa suatu matriks memiliki jumlah
      entri yang sama dengan dirinya sendiri.
      Untuk simetri, andaikan $A$ memiliki jumlah entri yang sama dengan~$B$;
      jelas bahwa $B$ kemudian memiliki jumlah entri yang sama dengan~$A$.
      Transitivitas juga mudah\Dash jika $A$ memiliki jumlah entri yang sama
      dengan $B$ dan $B$ memiliki jumlah entri yang sama dengan $C$, maka
      $A$ memiliki jumlah entri yang sama dengan $C$.
    \end{answer}
 \item \label{exer:INotJMakesRowOpsRev}
   Pembuktian \nearbylemma{le:RowOpsRev} menyebut syarat $i\neq j$ pada
   operasi kombinasi baris.
   \begin{exparts}
     \partsitem Tuliskan matriks $\nbyn{2}$ dengan entri-entri tak nol, lalu
        tunjukkan bahwa operasi $-1\cdot\rho_1+\rho_1$ tidak dibalik oleh
        $1\cdot\rho_1+\rho_1$.
     \partsitem Perluas pembuktian lema tersebut untuk menunjukkan secara
        eksplisit bagian yang menggunakan syarat $i\neq j$ pada kombinasi.
   \end{exparts}
   \begin{answer}
    \begin{exparts}
      \partsitem Sebagai contoh,
        \begin{equation*}
          \begin{mat}[r]
            1  &2  \\
            3  &4  
          \end{mat}
          \grstep{-\rho_1+\rho_1}
          \begin{mat}[r]
            0  &0  \\
            3  &4  
          \end{mat}
          \grstep{\rho_1+\rho_1}
          \begin{mat}[r]
            0  &0  \\
            3  &4  
          \end{mat}
        \end{equation*}
        membiarkan matriks dalam keadaan yang telah berubah.
      \partsitem Operasi berikut
        \begin{equation*}
          \begin{mat}
            \vdotswithin{a_{i,1}}                     \\
            a_{i,1}  &\cdots  &a_{i,n}  \\
            \vdotswithin{a_{i,1}}                     \\
            a_{j,1}  &\cdots  &a_{j,n}  \\
            \vdotswithin{a_{i,1}}                     
          \end{mat}
          \grstep{k\rho_i+\rho_j}
          \begin{mat}
            \vdotswithin{a_{i,1}}                                      \\
            a_{i,1}           &\cdots  &a_{i,n}          \\
            \vdotswithin{a_{i,1}}                                      \\
            ka_{i,1}+a_{j,1}  &\cdots  &ka_{i,n}+a_{j,n}  \\
            \vdotswithin{a_{i,1}}                     
          \end{mat}
        \end{equation*}      
        membiarkan baris ke-$i$ tetap karena pembatasan $i\neq j$.
        Karena baris ke-$i$ tidak berubah, operasi
        \begin{equation*}                                        
          \grstep{-k\rho_i+\rho_j}
          \begin{mat}
            \vdotswithin{a_{i,1}}                                      \\
            a_{i,1}           &\cdots  &a_{i,n}          \\
            \vdotswithin{a_{i,1}}                                      \\
            -ka_{i,1}+ka_{i,1}+a_{j,1}  &\cdots &-ka_{i,n}+ka_{i,n}+a_{j,n} \\
            \vdotswithin{a_{i,1}}                     
          \end{mat}
        \end{equation*}
        mengembalikan baris ke-$j$ ke keadaan semula.
        % (If $i=j$ then the third matrix would have entries of the 
        % form $-k(ka_{i,j}+a_{i,j})+ka_{i,j}+a_{i,j}$.)
    \end{exparts}
   \end{answer}
 \recommended \item \cite{Cleary}
  Tinjau himpunan siswa dalam suatu kelas.
  Manakah di antara relasi berikut yang merupakan relasi ekuivalensi?
  Jelaskan setiap jawaban dalam sedikitnya satu kalimat.
  \begin{exparts}
    \item Dua siswa $x,y$ berelasi jika $x$ telah mengikuti setidaknya
      sebanyak mata pelajaran matematika yang telah diikuti $y$.
    \item Siswa $x,y$ berelasi jika nama mereka diawali huruf yang sama.
  \end{exparts}
  \begin{answer}
    Agar menjadi relasi ekuivalensi, setiap relasi harus refleksif, simetris,
    dan transitif.
    \begin{exparts}
      \item Relasi ini tidak simetris karena jika $x$ telah mengikuti
        $4$~mata pelajaran dan $y$ telah mengikuti $3$, maka $x$ berelasi
        dengan $y$, tetapi $y$ tidak berelasi dengan $x$.
      \item Relasi ini refleksif karena nama $x$ diawali huruf yang sama dengan
        nama $x$ sendiri.
        Relasi ini simetris karena jika nama $x$ diawali huruf yang sama dengan
        nama $y$, maka nama $y$ juga diawali huruf yang sama dengan nama~$x$.
        Relasi ini transitif karena jika nama $x$ diawali huruf yang sama dengan
        nama~$y$ dan nama $y$ diawali huruf yang sama dengan nama $z$, maka
        nama $x$ diawali huruf yang sama dengan nama $z$.
        Jadi, relasi ini merupakan relasi ekuivalensi.
    \end{exparts}
  \end{answer}
  \item
  Tunjukkan bahwa masing-masing relasi berikut merupakan relasi ekuivalensi
  pada himpunan matriks $\nbyn{2}$.
  Jelaskan kelas-kelas ekuivalensinya.
  \begin{exparts}
    \item Dua matriks berelasi jika hasil kali entri pada diagonalnya sama,
      yakni jika hasil kali entri kiri atas dan kanan bawahnya sama.
    \item Dua matriks berelasi jika keduanya memiliki sedikitnya satu entri
      bernilai~$1$, atau jika keduanya sama-sama tidak memilikinya.
  \end{exparts}
  \begin{answer}
    Untuk masing-masing relasi, kita harus memeriksa ketiga syarat
    refleksivitas, simetri, dan transitivitas.
    \begin{exparts}
      \item Setiap matriks jelas memiliki hasil kali diagonal yang sama dengan
        dirinya sendiri, sehingga relasi tersebut refleksif.
        Relasi itu simetris karena jika $A$ memiliki hasil kali diagonal yang
        sama dengan~$B$, yaitu jika
        $a_{1,1}\cdot a_{2,2}=b_{1,1}\cdot b_{2,2}$, maka $B$ memiliki hasil
        kali yang sama dengan~$A$.
 
        Transitivitas serupa: andaikan hasil kali $A$ adalah~$r$ dan sama
        dengan hasil kali $B$.
        Andaikan pula hasil kali $B$ sama dengan hasil kali $C$.
        Maka ketiganya memiliki hasil kali~$r$, sehingga hasil kali $A$ sama
        dengan hasil kali $C$.

        Terdapat satu kelas ekuivalensi untuk setiap bilangan riil, yaitu kelas
        yang memuat semua matriks $\nbyn{2}$ dengan hasil kali diagonal sama
        dengan bilangan riil tersebut.
      \item Untuk refleksivitas, jika matriks $A$ memiliki entri~$1$, maka ia
        berelasi dengan dirinya sendiri; jika tidak, ia juga tetap berelasi
        dengan dirinya sendiri.
        Simetri juga memiliki dua kasus: andaikan $A$ dan~$B$ berelasi.
        Jika $A$ memiliki entri~$1$, maka $B$ juga memilikinya, sehingga $B$
        berelasi dengan~$A$.
        Jika $A$ tidak memiliki entri~$1$, maka $B$ juga tidak, sehingga sekali
        lagi $B$ berelasi dengan $A$.
 
        Untuk transitivitas, andaikan $A$ berelasi dengan~$B$ dan $B$ berelasi
        dengan~$C$.
        Jika $A$ memiliki entri~$1$, maka $B$ juga memilikinya; karena $B$
        berelasi dengan~$C$, $C$ juga memilikinya. Jadi, $A$ berelasi
        dengan~$C$.
        Demikian pula, jika $A$ tidak memiliki entri~$1$, maka $B$ juga tidak,
        dan akibatnya $C$ juga tidak. Jadi, $A$ berelasi dengan~$C$.

        Tepat ada dua kelas ekuivalensi: satu memuat semua matriks $\nbyn{2}$
        yang memiliki sedikitnya satu entri bernilai~$1$, dan yang lain memuat
        semua matriks yang tidak memiliki entri bernilai~$1$.
    \end{exparts}
  \end{answer}
  \item Tunjukkan bahwa masing-masing relasi berikut bukan relasi ekuivalensi
    pada himpunan matriks $\nbyn{2}$.
    \begin{exparts}
      \item Dua matriks $A,B$ berelasi jika $a_{1,1}=-b_{1,1}$.
      \item Dua matriks berelasi jika selisih jumlah entrinya kurang dari
        $5$, yakni $A$ berelasi dengan~$B$ jika
        $\absval{(a_{1,1}+\cdots+a_{2,2})-(b_{1,1}+\cdots+b_{2,2})}<5$.
    \end{exparts}
    \begin{answer}
      \begin{exparts}
        \item Relasi ini tidak refleksif.
          Sebagai contoh, setiap matriks dengan entri kiri atas~$1$ tidak
          berelasi dengan dirinya sendiri.
        \item Relasi ini tidak transitif.
          Untuk ketiga matriks berikut, $A$ berelasi dengan~$B$ dan $B$
          berelasi dengan~$C$, tetapi $A$ tidak berelasi dengan~$C$.
          \begin{equation*}
            A=\begin{mat}
              0 &0 \\
              0 &0
            \end{mat},\quad
            B=\begin{mat}
              4 &0 \\
              0 &0
            \end{mat},\quad
            C=\begin{mat}
              8 &0 \\
              0 &0
            \end{mat},\quad
          \end{equation*}
      \end{exparts}
    \end{answer}
\end{exercises}




















\subsection{Lema Kombinasi Linear}
Kita akan menutup bab ini dengan membuktikan bahwa setiap matriks ekuivalen
baris dengan tepat satu matriks berbentuk eselon tereduksi.
Gagasan-gagasan di sini akan muncul kembali dan dikembangkan lebih lanjut
dalam bab berikutnya.
% Here is an informal argument that the reduced 
% echelon form version of a matrix is unique.
% Consider again the example that started this section of a matrix that
% reduces to three different echelon form matrices.
% The first matrix of the three is the natural echelon form version.
% The second matrix is the same as 
% the first except that a row has been halved.
% The third matrix, too, is just a cosmetic variant of the first. 
% The definition of reduced echelon form outlaws this kind of fooling around.
% In reduced echelon form,
% halving a row is not possible because that would
% change the row's leading entry away from one, and
% neither is combining rows possible, because then a leading entry would no
% longer be alone in its column.

% This informal justification is not a proof;
% the argument shows that no two different reduced echelon form matrices
% are related by a single row operation step, but the argument does not
% rule out the possibility that two different reduced echelon form
% matrices could be related by multiple steps.
% Before we go to the proof, we finish this subsection by 
% rephrasing our work in a terminology that will be enlightening.
Pengamatan kuncinya menyangkut cara operasi baris mengubah suatu matriks
menjadi matriks lain: baris-baris baru merupakan kombinasi linear dari
baris-baris lama.

\begin{example}
Tinjau reduksi Gauss--Jordan berikut.
\begin{align*}
  \begin{amat}{2}
    2  &1  &0  \\
    1  &3  &5
  \end{amat}
  &\grstep{-(1/2)\rho_1+\rho_2}
  \begin{amat}{2}
    2  &1    &0  \\
    0  &5/2  &5
  \end{amat}                           \\
  &\grstep[(2/5)\rho_2]{(1/2)\rho_1}
  \begin{amat}{2}
    1  &1/2  &0  \\
    0  &1    &2
  \end{amat}                        \\
  &\grstep{-(1/2)\rho_2+\rho_1}\;
  \begin{amat}{2}
    1  &0    &-1  \\
    0  &1    &2
  \end{amat}                       
\end{align*}
Dengan menamai matriks-matriks tersebut
$A\rightarrow D\rightarrow G\rightarrow B$ dan menuliskan baris-baris $A$
sebagai $\alpha_1$ dan $\alpha_2$, dan seterusnya, kita memperoleh
\begin{align*}  % multline looks worse
  \left(\begin{array}{l}
    \alpha_1 \\
    \alpha_2
  \end{array}\right)
  &\grstep{-(1/2)\rho_1+\rho_2}
  \left(\begin{array}{l}
    \delta_1=\alpha_1 \\
    \delta_2=-(1/2)\alpha_1+\alpha_2
  \end{array}\right)                          \\
  &\grstep[(2/5)\rho_2]{(1/2)\rho_1}
  \left(\begin{array}{l}
    \gamma_1=(1/2)\alpha_1 \\
    \gamma_2=-(1/5)\alpha_1+(2/5)\alpha_2
  \end{array}\right)                          \\
  &\grstep{-(1/2)\rho_2+\rho_1}
  \left(\begin{array}{l}
    \beta_1=(3/5)\alpha_1-(1/5)\alpha_2  \\
    \beta_2=-(1/5)\alpha_1+(2/5)\alpha_2
  \end{array}\right)                         
\end{align*} 
\end{example}

\begin{example}
Fakta bahwa operasi Gauss mengombinasikan baris secara linear juga berlaku
jika terjadi pertukaran baris.
Dengan \( A \), \( D \), \( G \), dan \( B \) berikut,
\begin{equation*}
    \begin{mat}[r]
       0  &2  \\
       1  &1
     \end{mat}
    \grstep{\rho_1\leftrightarrow\rho_2}
    \begin{mat}[r]
       1  &1  \\
       0  &2
     \end{mat}                  
    \grstep{(1/2)\rho_2}
    \begin{mat}[r]
       1  &1  \\
       0  &1
     \end{mat}                   
    \grstep{-\rho_2+\rho_1}
    \begin{mat}[r]
       1  &0  \\
       0  &1
     \end{mat}
\end{equation*}
kita memperoleh hubungan linear berikut.
\begin{multline*}
  \left(\begin{array}{l}
    \vec{\alpha}_1 \\
    \vec{\alpha}_2
  \end{array}\right)
  \,\grstep{\rho_1\leftrightarrow\rho_2}\,
  \left(\begin{array}{l}
     \vec{\delta}_1 =\vec{\alpha}_2   \\
     \vec{\delta}_2 =\vec{\alpha}_1
  \end{array}\right)                                                 
  \,\grstep{(1/2)\rho_2}\,
  \left(\begin{array}{l}
     \vec{\gamma}_1 =\vec{\alpha}_2  \\
     \vec{\gamma}_2 =(1/2)\vec{\alpha}_1
  \end{array}\right)                                               \\
  \,\grstep{-\rho_2+\rho_1}\,
  \left(\begin{array}{l}
     \vec{\beta}_1 =(-1/2)\vec{\alpha}_1+1\cdot\vec{\alpha}_2   \\
     \vec{\beta}_2 =(1/2)\vec{\alpha}_1
  \end{array}\right)
\end{multline*}
\end{example}

Ringkasnya, Metode Gauss secara sistematis menentukan urutan kombinasi linear
baris yang sesuai.

% \begin{definition}
% A \definend{linear combination}\index{linear combination} of 
% \( x_1,\ldots,x_m \)
% is an expression of the form $c_1x_1+c_2x_2+\,\cdots\,+c_mx_m$
% where the \( c \)'s are scalars.
%\end{definition}

% \noindent (We have already used the phrase 
% `linear combination' in this book.
% The meaning is unchanged, but the next result's statement makes
% a more formal definition in order.)

\begin{lemma}[Lema Kombinasi Linear] \label{lm:LinearCombinationLemma} \index{Lema Kombinasi Linear}
%<*lm:LinearCombinationLemma>
Kombinasi linear dari kombinasi-kombinasi linear merupakan kombinasi linear.
%</lm:LinearCombinationLemma>
\end{lemma}

\begin{proof}
%<*pf:LinearCombinationLemma>
Diberikan kumpulan kombinasi linear dari variabel-variabel
$x_1,\ldots,x_n$,
$c_{1,1}x_1+\dots+c_{1,n}x_n$ hingga
$c_{m,1}x_1+\dots+c_{m,n}x_n$, tinjau kombinasi dari semuanya,
\begin{equation*}
  d_1(c_{1,1}x_1+\dots+c_{1,n}x_n)\,+\dots+\,d_m(c_{m,1}x_1+\dots+c_{m,n}x_n)
\end{equation*}
dengan semua $d_i$ dan $c_{i,j}$ merupakan skalar.
Mendistribusikan $d$ tersebut dan mengelompokkan ulang menghasilkan
\begin{equation*}
  %&=d_1c_{1,1}x_1+\dots+d_1c_{1,n}x_n\,
  % +d_2c_{2,1}x_1+\dots+\,
  % d_mc_{1,1}x_1+\dots+d_mc_{1,n}x_n         \\
  =(d_1c_{1,1}+\dots+d_mc_{m,1})x_1\,+\dots+\,(d_1c_{1,n}+\dots+d_mc_{m,n})x_n
\end{equation*}
yang juga merupakan kombinasi linear dari $x_1,\ldots,x_n$.
%</pf:LinearCombinationLemma>
\end{proof}


\begin{corollary} \label{cor:RowsOfEqMatsLinCombos}
%<*co:RowsOfEqMatsLinCombos>
Jika suatu matriks mereduksi menjadi matriks lain, setiap baris matriks kedua
merupakan kombinasi linear dari baris-baris matriks pertama.
%</co:RowsOfEqMatsLinCombos>
\end{corollary}

% The proof 
% uses induction. % \appendrefs{mathematical induction}\spacefactor=1000 %
% Before we proceed, here is an outline of the argument.
% For the base step, we
% will verify that the proposition is true when reduction 
% can be done in zero row operations.
% For the inductive step, we will 
% argue that if being able to reduce the first matrix to the second in some
% number $t\geq 0$ of operations implies that each row of the second is a linear
% combination of the rows of the first, then being able to reduce the first to
% the second in $t+1$ operations implies the same thing.
% Together these prove the result because  
% the base step shows that it is true in the zero operations case,
% and then the inductive step
% implies that it is true in the one operation case, and then the inductive step
% applied again gives that it is therefore true for two operations, etc.

\begin{proof}
%<*pf:RowsOfEqMatsLinCombos0>
Untuk sebarang dua matriks $A$ dan~$B$ yang dapat saling direduksi, terdapat
jumlah minimum operasi baris yang mengubah salah satunya menjadi yang lain.
Kita menggunakan induksi pada jumlah tersebut.

Pada langkah dasar, jika kita dapat berpindah dari satu matriks ke matriks lain
dengan nol operasi reduksi, kedua matriks itu sama.
Dengan demikian, setiap baris $B$ secara langsung merupakan kombinasi dari
baris-baris $A$,
$\vec{\beta}_i
  =0\cdot\vec{\alpha}_1+\cdots+1\cdot\vec{\alpha}_i+\cdots+0\cdot\vec{\alpha}_m$.
%</pf:RowsOfEqMatsLinCombos0>

%<*pf:RowsOfEqMatsLinCombos1>
Untuk langkah induksi, andaikan hipotesis induksi berikut:~untuk $k\geq 0$,
setiap matriks yang dapat diperoleh dari \( A \) dalam \( k \) operasi atau
kurang memiliki baris-baris yang merupakan kombinasi linear dari baris-baris
$A$.
Tinjau matriks~$B$ sedemikian sehingga reduksi \( A \) menjadi~\( B \)
memerlukan $k+1$ operasi.
Dalam reduksi tersebut terdapat matriks kedua dari akhir,~$G$, sehingga
$A\longrightarrow\cdots\longrightarrow G\longrightarrow B$.
Hipotesis induksi berlaku bagi \( G \) karena jaraknya hanya $k$ langkah dari
\( A \).
Artinya, setiap baris \( G \) merupakan kombinasi linear dari baris-baris
\( A \).
%</pf:RowsOfEqMatsLinCombos1>

%<*pf:RowsOfEqMatsLinCombos2>
Kita akan memastikan bahwa baris-baris~$B$ merupakan kombinasi linear dari
baris-baris~$G$.
Lema Kombinasi Linear, \nearbylemma{lm:LinearCombinationLemma}, kemudian
menunjukkan bahwa baris-baris~$B$ merupakan kombinasi linear dari baris-baris
~$A$.

Jika operasi baris yang mengubah \( G \) menjadi~\( B \) adalah pertukaran,
baris-baris $B$ hanyalah baris-baris $G$ dalam urutan berbeda dan setiap baris
$B$ merupakan kombinasi linear dari baris-baris $G$.
Jika operasi yang mengubah $G$ menjadi~$B$ adalah perkalian suatu baris dengan
skalar~$c\rho_i$, maka $\vec{\beta}_i=c\vec{\gamma}_i$ dan baris-baris lainnya
tidak berubah.
Terakhir, jika operasi barisnya menambahkan kelipatan suatu baris ke baris lain,
$r\rho_i+\rho_j$, hanya baris~$j$ dari $B$ yang berbeda dari baris bersesuaian
dalam~$G$, dan
$\vec{\beta}_j=r\vec{\gamma}_i+\vec{\gamma}_j$, yang memang merupakan
kombinasi linear dari baris-baris $G$.
%</pf:RowsOfEqMatsLinCombos2>

Karena kita telah membuktikan langkah dasar maupun langkah induksi, proposisi
tersebut mengikuti prinsip induksi matematika.
\end{proof}

Sekarang kita memahami bahwa Metode Gauss membangun kombinasi linear dari
baris-baris.
Namun, tentu saja tujuannya adalah mencapai bentuk eselon karena bentuk itu
merupakan versi sistem linear yang sangat mendasar, dengan variabel-variabel
yang telah diisolasi.
Sebagai contoh, dalam matriks
\begin{equation*}
  R=\begin{mat}[r]
    2  &3  &7  &8  &0  &0  \\
    0  &0  &1  &5  &1  &1  \\
    0  &0  &0  &3  &3  &0  \\
    0  &0  &0  &0  &2  &1
  \end{mat}
\end{equation*}
$x_1$ telah dieliminasi dari persamaan $x_5$.
Artinya, Metode Gauss telah membuat baris $x_5$ bebas, dalam suatu arti,
dari baris $x_1$.

Hasil berikut menjadikan intuisi ini tepat.
Kita kadang menyebut Metode Gauss sebagai eliminasi Gauss.
Yang dieliminasinya adalah hubungan linear antarbaris.

\begin{lemma}      \label{le:EchFormNoLinCombo}
%<*lm:EchFormNoLinCombo>
Dalam matriks berbentuk eselon, tidak ada baris tak nol yang merupakan
kombinasi linear dari baris-baris tak nol lainnya.
%</lm:EchFormNoLinCombo>
\end{lemma}

\begin{proof}
%<*pf:EchFormNoLinCombo0>
Misalkan $R$ matriks berbentuk eselon dan tinjau baris-baris tak nolnya.
Pertama, perhatikan bahwa jika suatu baris dituliskan sebagai kombinasi
baris-baris lainnya,
$\vec{\rho}_i=c_1\vec{\rho}_1+\cdots+c_{i-1}\vec{\rho}_{i-1}+
               c_{i+1}\vec{\rho}_{i+1}+\cdots+c_m\vec{\rho}_m$
maka persamaan itu dapat ditulis ulang sebagai
\begin{equation*}
   \vec{0}=c_1\vec{\rho}_1+\cdots+c_{i-1}\vec{\rho}_{i-1}+c_i\vec{\rho}_i+
               c_{i+1}\vec{\rho}_{i+1}+\cdots+c_m\vec{\rho}_m
  \tag{$*$}
\end{equation*}
dengan tidak semua koefisien bernilai nol; khususnya, $c_i=-1$.
Pernyataan sebaliknya juga berlaku:~diberikan persamaan~($*$) dengan suatu
$c_i\neq 0$, kita dapat menyatakan $\vec{\rho}_i$ sebagai kombinasi baris-baris
lainnya dengan memindahkan $c_i\vec{\rho}_i$ ke kiri dan membagi dengan
$-c_i$.
Jadi, untuk membuktikan lema tersebut, cukup kita tunjukkan bahwa semua
koefisien dalam~($*$) bernilai~$0$.
Untuk itu, kita menggunakan induksi pada nomor baris~$i$.
%</pf:EchFormNoLinCombo0>

%<*pf:EchFormNoLinCombo1>
Kasus dasarnya adalah baris pertama,~$i=1$ (jika tidak ada baris tak nol
seperti itu, sehingga $R$ merupakan matriks nol, lema berlaku secara hampa).
Misalkan $\ell_i$ nomor kolom entri utama pada baris~$i$.
Tinjau entri setiap baris yang berada dalam kolom~$\ell_1$.
Persamaan~($*$) memberikan
\begin{equation*}
  0=c_1r_{1,\ell_1}+c_2r_{2,\ell_1}+\cdots+c_mr_{m,\ell_1}
  \tag{$**$}
\end{equation*}
Matriks berada dalam bentuk eselon sehingga setiap baris setelah baris pertama
memiliki entri nol dalam kolom itu,
$r_{2,\ell_1}=\cdots=r_{m,\ell_1}=0$.
Jadi, persamaan~($**$) menunjukkan bahwa $c_1=0$ karena
$r_{1,\ell_1}\neq 0$ sebagai entri utama baris tersebut.
%</pf:EchFormNoLinCombo1>

%<*pf:EchFormNoLinCombo2>
Langkah induksi hampir sama dengan langkah dasar.
Tinjau kembali persamaan~($*$).
Kita akan membuktikan bahwa jika koefisien $c_i$ bernilai $0$ untuk setiap
indeks baris $i\in\set{1,\ldots,k}$, maka $c_{k+1}$ juga bernilai $0$.
Kita memusatkan perhatian pada entri-entri kolom~$\ell_{k+1}$.
\begin{equation*}
  0=c_1r_{1,\ell_{k+1}}+\cdots+c_{k+1}r_{k+1,\ell_{k+1}}+\cdots+c_mr_{m,\ell_{k+1}}
\end{equation*}
Menurut hipotesis induksi, $c_1$, \ldots, $c_k$ semuanya bernilai $0$, sehingga
persamaan tersebut tereduksi menjadi
$0=c_{k+1}r_{k+1,\ell_{k+1}}+\cdots+c_mr_{m,\ell_{k+1}}$.
Matriks berada dalam bentuk eselon sehingga entri
$r_{k+2,\ell_{k+1}}$, \ldots, $r_{m,\ell_{k+1}}$ semuanya bernilai~$0$.
Jadi, $c_{k+1}=0$ karena $r_{k+1,\ell_{k+1}}\neq 0$ sebagai entri utama.
%</pf:EchFormNoLinCombo2>
\end{proof}

Dengan hasil itu, kita siap menunjukkan bahwa hasil akhir reduksi
Gauss--Jordan bersifat unik.

\begin{theorem}
\label{th:ReducedEchelonFormIsUnique}
%<*th:ReducedEchelonFormIsUnique>
Setiap matriks ekuivalen baris dengan tepat satu matriks berbentuk eselon
tereduksi.
%</th:ReducedEchelonFormIsUnique>
\end{theorem}

\begin{proof} \cite{Yuster}
%<*pf:ReducedEchelonFormIsUnique0>
Tetapkan jumlah baris \( m \).
Kita menggunakan induksi pada jumlah kolom \( n \).

Kasus dasarnya adalah matriks dengan \( n=1 \) kolom.
Jika matriks itu matriks nol, bentuk eselonnya adalah matriks nol.
Jika sebaliknya matriks memiliki suatu entri tak nol, ketika dibawa ke bentuk
eselon tereduksi, matriks itu harus memiliki sedikitnya satu entri tak nol,
yang harus berupa \( 1 \) pada baris pertama.
Dalam kedua kasus, bentuk eselon tereduksinya unik.
%</pf:ReducedEchelonFormIsUnique0>

%<*pf:ReducedEchelonFormIsUnique1>
Untuk langkah induksi, andaikan \( n>1 \) dan semua matriks dengan \( m \)
baris dan kurang dari~\( n \) kolom memiliki bentuk eselon tereduksi yang unik.
Tinjau matriks \( \nbym{m}{n} \), \( A \), dan andaikan \( B \) serta
\( C \) merupakan dua matriks berbentuk eselon tereduksi yang diperoleh dari
\( A \).
Kita akan menunjukkan bahwa keduanya harus sama.
%</pf:ReducedEchelonFormIsUnique1>

%<*pf:ReducedEchelonFormIsUnique2>
Misalkan \( \hat{A} \) matriks yang terdiri atas \( n-1 \) kolom pertama dari
\( A \).
Perhatikan bahwa
% if an $n$~column matrix is in reduced echelon form then any initial
% set of its columns is also in reduced echelon form, so \( \hat{A} \)
% is in reduced echelon form. 
setiap urutan operasi baris yang membawa \( A \) ke bentuk eselon tereduksi
juga membawa \( \hat{A} \) ke bentuk eselon tereduksi.
Menurut hipotesis induksi, bentuk eselon tereduksi \( \hat{A} \) ini unik;
jadi, jika \( B \) dan \( C \) berbeda, perbedaannya harus terdapat dalam
kolom~\( n \).
%</pf:ReducedEchelonFormIsUnique2>

Kita menyelesaikan langkah induksi dan pembuktian dengan menunjukkan bahwa
keduanya tidak mungkin hanya berbeda dalam kolom tersebut.
%<*pf:ReducedEchelonFormIsUnique3>
Tinjau sistem persamaan homogen dengan \( A \) sebagai matriks koefisien.
\begin{equation*}
  \begin{linsys}{4}
    a_{1,1}x_1  &+  &a_{1,2}x_2  &+  &\cdots  &+  &a_{1,n}x_n  &=  &0  \\
    a_{2,1}x_1  &+  &a_{2,2}x_2  &+  &\cdots  &+  &a_{2,n}x_n  &=  &0  \\
              &&&&&&&\vdotswithin{=}  \\
    a_{m,1}x_1  &+  &a_{m,2}x_2  &+  &\cdots  &+  &a_{m,n}x_n  &=  &0  
  \end{linsys}
  \tag{$*$}
\end{equation*}
Menurut Teorema~I.\ref{th:GaussMethod} di Bab Satu,
%the first theorem of this chapter 
himpunan solusi sistem tersebut sama dengan himpunan solusi sistem milik $B$,
\begin{equation*}
  \begin{linsys}{4}
    b_{1,1}x_1  &+  &b_{1,2}x_2  &+  &\cdots  &+  &b_{1,n}x_n  &=  &0  \\
    b_{2,1}x_1  &+  &b_{2,2}x_2  &+  &\cdots  &+  &b_{2,n}x_n  &=  &0  \\
               &&&&&&&\vdotswithin{=}  \\
    b_{m,1}x_1  &+  &b_{m,2}x_2  &+  &\cdots  &+  &b_{m,n}x_n  &=  &0  
  \end{linsys}
  \tag{$**$}
\end{equation*}
dan juga sama dengan himpunan solusi sistem milik $C$.
\begin{equation*}
  \quad
  \begin{linsys}{4}
    c_{1,1}x_1  &+  &c_{1,2}x_2  &+  &\cdots  &+  &c_{1,n}x_n  &=  &0  \\
    c_{2,1}x_1  &+  &c_{2,2}x_2  &+  &\cdots  &+  &c_{2,n}x_n  &=  &0  \\
               &&&&&&&\vdotswithin{=}  \\
    c_{m,1}x_1  &+  &c_{m,2}x_2  &+  &\cdots  &+  &c_{m,n}x_n  &=  &0  
  \end{linsys}
  \tag{$\mathord{*}\mathord{*}\mathord{*}$}
\end{equation*}
%</pf:ReducedEchelonFormIsUnique3>
%<*pf:ReducedEchelonFormIsUnique4>
Karena \( B \) dan \( C \) hanya berbeda dalam kolom~\( n \), andaikan
keduanya berbeda pada baris~\( i \).
Kurangkan baris ke-\( i \) pada ($\mathord{*}\mathord{*}\mathord{*}$) dari
baris ke-\( i \) pada ($**$) untuk memperoleh persamaan
\( (b_{i,n}-c_{i,n})\cdot x_n=0 \).
Kita telah mengandaikan \( b_{i,n}\neq c_{i,n} \), sehingga diperoleh
$x_n=0$.
Jadi, $x_n$ bukan variabel bebas; karena itu, dalam ($**$) dan
~($\mathord{*}\mathord{*}\mathord{*}$), kolom ke-\( n \) memuat entri utama
suatu baris, sebab dalam matriks berbentuk eselon, setiap kolom tanpa entri
utama berkaitan dengan variabel bebas.

Namun, karena \( B \) dan \( C \) sama pada \( n-1 \) kolom pertama, menurut
definisi bentuk eselon tereduksi, entri utama keduanya dalam kolom ke-\( n \)
berada pada baris yang sama.
Kedua entri utama itu harus bernilai $1$ dan harus menjadi satu-satunya entri
tak nol dalam kolom tersebut.
Jadi, \( B=C \).
%</pf:ReducedEchelonFormIsUnique4>
\end{proof}

Kita telah bertanya apakah sebarang dua versi berbentuk eselon dari suatu sistem
linear memiliki jumlah variabel bebas yang sama dan, jika ya, apakah variabel
bebasnya persis sama.
Dengan hasil sebelumnya, kita dapat menjawab ``ya'' untuk kedua pertanyaan.
Tidak ada sistem linear yang, misalnya, menghasilkan $y$ dan $z$ sebagai
variabel bebas ketika Metode Gauss diterapkan dengan satu cara, tetapi
menghasilkan $y$ dan $w$ sebagai variabel bebas ketika diterapkan dengan cara
lain.

Sebelum pembuktian, ingat kembali perbedaan antara variabel bebas dan
parameter.
Sistem berikut
\begin{equation*}
  \begin{linsys}{3}
    x &+ &y &  &  &= &1 \\ 
      &  &y &+ &z &= &2  
  \end{linsys}
\end{equation*}
memiliki satu variabel bebas,~$z$, karena hanya variabel itu yang bukan
variabel utama suatu baris.
Kita biasanya melakukan parametrisasi menggunakan variabel bebas tersebut,
$y=2-z$, $x=-1+z$, tetapi kita juga dapat melakukan parametrisasi menggunakan
variabel lain, misalnya $z=2-y$, $x=1-y$.
Jadi, himpunan parameter tidak unik; himpunan variabel bebaslah yang unik.

\begin{corollary}
Jika dari suatu sistem linear awal kita memperoleh dua sistem berbentuk eselon
yang berbeda melalui Metode Gauss, keduanya memiliki variabel bebas yang sama.
\end{corollary}

\begin{proof}
Hasil sebelumnya menyatakan bahwa bentuk eselon tereduksi bersifat unik.
Kita berpindah dari versi berbentuk eselon mana pun ke bentuk eselon tereduksi
dengan mengeliminasi ke arah atas; jadi, setiap versi berbentuk eselon suatu
sistem memiliki variabel bebas yang sama dengan versi bentuk eselon
tereduksinya.
\end{proof}

Kita tutup dengan rangkuman.
Dalam Metode Gauss, kita mulai dengan suatu matriks, lalu memperoleh urutan
matriks-matriks lain.
Kita mendefinisikan dua matriks berelasi jika salah satunya dapat diperoleh
dari yang lain.
Relasi tersebut merupakan relasi ekuivalensi, %\appendrefs{equivalence relation} 
yang disebut ekuivalensi baris, sehingga mempartisi himpunan semua matriks
menjadi kelas-kelas ekuivalensi baris.
\begin{center}
  \includegraphics{gr/mp/ch1.30}
\end{center}
(Kelas yang digambarkan memuat tak hingga banyak matriks, tetapi ruang yang
tersedia hanya cukup untuk menunjukkan dua.)
Kita telah membuktikan bahwa terdapat tepat satu matriks berbentuk eselon
tereduksi dalam setiap kelas ekuivalensi baris.
%<*ReducedEchelonFormIsCanonicalForm>
Jadi, bentuk eselon tereduksi merupakan
bentuk kanonis\appendrefs{wakil kanonis}\spacefactor=1000
\index{bentuk kanonis!untuk ekuivalensi baris}\index{wakil!untuk kelas ekuivalensi baris}
bagi ekuivalensi baris: matriks-matriks berbentuk eselon tereduksi merupakan
wakil kelas-kelas tersebut.
%</ReducedEchelonFormIsCanonicalForm>
\begin{center}
  \includegraphics{gr/mp/ch1.31}
\end{center}

Gagasannya adalah bahwa salah satu cara memahami situasi matematika ialah
dengan mengklasifikasikan kasus-kasus yang dapat terjadi.
Hal ini merupakan tema buku ini dan telah kita jumpai beberapa kali.
Kita mengklasifikasikan himpunan solusi sistem linear menjadi kasus tanpa
anggota, beranggota satu, dan beranggota tak hingga.
Kita juga mengklasifikasikan sistem linear dengan jumlah persamaan sama dengan
jumlah variabel tak diketahui menjadi kasus nonsingular dan singular.

Di sini, ketika menyelidiki ekuivalensi baris, kita mengetahui bahwa himpunan
semua matriks terbagi menjadi kelas-kelas ekuivalensi baris dan sekarang kita
memiliki cara untuk mengidentifikasi setiap kelas\Dash kita dapat
memandang matriks-matriks dalam suatu kelas sebagai hasil operasi baris dari
matriks berbentuk eselon tereduksi unik dalam kelas itu.

Dalam istilah yang lebih operasional, keunikan bentuk eselon tereduksi
memungkinkan kita menjawab pertanyaan tentang kelas dengan mengubahnya menjadi
pertanyaan tentang wakilnya.
Sebagai contoh, seperti dijanjikan pada awal bagian ini, sekarang kita dapat
menentukan apakah suatu matriks dapat diperoleh dari matriks lain melalui
reduksi baris.
Kita menerapkan prosedur Gauss--Jordan pada keduanya dan memeriksa apakah
keduanya menghasilkan bentuk eselon tereduksi yang sama.

\begin{example}  \label{ex:MatsNotRowEq}
Matriks-matriks berikut tidak ekuivalen baris
\begin{equation*}
  \begin{mat}[r]
    1  &-3  \\
   -2  &6
  \end{mat}
  \qquad
  \begin{mat}[r]
    1  &-3  \\
   -2  &5
  \end{mat}
\end{equation*}
karena bentuk eselon tereduksinya tidak sama.
\begin{equation*}
  \begin{mat}[r]
    1  &-3  \\
    0  &0
  \end{mat}
  \qquad
  \begin{mat}[r]
    1  &0   \\
    0  &1
  \end{mat}
\end{equation*}
\end{example}

\begin{example}
Setiap matriks nonsingular \( \nbyn{3} \) tereduksi melalui Gauss--Jordan
menjadi matriks berikut.
\begin{equation*}
    \begin{mat}[r]
      1  &0  &0 \\
      0  &1  &0 \\
      0  &0  &1
    \end{mat}
\end{equation*}
\end{example}

\begin{example} \label{ex:RowEqClassTwoTwoMats}
Kita dapat mendeskripsikan semua kelas dengan mendaftarkan semua matriks
berbentuk eselon tereduksi yang mungkin.
Setiap matriks $\nbyn{2}$ berada dalam salah satu kelas berikut:~kelas matriks
yang ekuivalen baris dengan
\begin{equation*}
  \begin{mat}[r]
     0  &0  \\
     0  &0
  \end{mat}
\end{equation*}
tak hingga banyak kelas matriks yang ekuivalen baris dengan salah satu matriks
berjenis
\begin{equation*}
  \begin{mat}
     1  &a  \\
     0  &0
  \end{mat}
\end{equation*}
dengan \( a\in\Re \) (termasuk $a=0$), kelas matriks yang ekuivalen baris
dengan
\begin{equation*}
  \begin{mat}[r]
     0  &1  \\
     0  &0
  \end{mat}
\end{equation*}
dan kelas matriks yang ekuivalen baris dengan
\begin{equation*}
  \begin{mat}[r]
     1  &0  \\
     0  &1
  \end{mat}
\end{equation*}
(ini adalah kelas matriks nonsingular $\nbyn{2}$).
\end{example}



\begin{exercises}
  \recommended \item 
    Tentukan apakah pasangan-pasangan matriks berikut ekuivalen baris.
    \begin{exparts*}
       \partsitem \(
           \begin{mat}[r]
             1  &2  \\
             4  &8
           \end{mat}, 
           \begin{mat}[r]
             0  &1  \\
             1  &2
           \end{mat} \)
       \partsitem \(
           \begin{mat}[r]
             1  &0  &2  \\
             3  &-1 &1  \\
             5  &-1 &5
           \end{mat},  
           \begin{mat}[r]
             1  &0  &2  \\
             0  &2  &10 \\
             2  &0  &4
           \end{mat} \)
       \partsitem \(
           \begin{mat}[r]
             2  &1  &-1 \\
             1  &1  &0  \\
             4  &3  &-1
           \end{mat},  
           \begin{mat}[r]
             1  &0  &2  \\
             0  &2  &10 \\
           \end{mat} \)
       \partsitem \(
           \begin{mat}[r]
             1  &1  &1  \\
            -1  &2  &2
           \end{mat},  
           \begin{mat}[r]
             0  &3  &-1 \\
             2  &2  &5
           \end{mat} \)
       \partsitem \(
           \begin{mat}[r]
             1  &1  &1  \\
             0  &0  &3
           \end{mat},  
           \begin{mat}[r]
             0  &1  &2  \\
             1  &-1 &1
           \end{mat} \)
    \end{exparts*}
    \begin{answer}
      Bawa setiap matriks ke bentuk eselon tereduksi, lalu bandingkan.
      \begin{exparts}
        \partsitem Matriks pertama menghasilkan
          \begin{equation*}
            \grstep{-4\rho_1+\rho_2}
            \begin{mat}[r]
              1  &2  \\
              0  &0
            \end{mat}
          \end{equation*}
          sedangkan matriks kedua menghasilkan
          \begin{equation*}
            \grstep{\rho_1\leftrightarrow\rho_2}
            \begin{mat}[r]
              1  &2  \\
              0  &1
            \end{mat}
            \grstep{-2\rho_2+\rho_1}
            \begin{mat}[r]
              1  &0  \\
              0  &1
            \end{mat}
          \end{equation*}
          Kedua matriks berbentuk eselon tereduksi itu tidak identik, sehingga
          matriks semula tidak ekuivalen baris.
        \partsitem Matriks pertama menjadi
          \begin{equation*}
            \grstep[-5\rho_1+\rho_3]{-3\rho_1+\rho_2}
            \begin{mat}[r]
              1  &0  &2  \\
              0  &-1 &-5 \\
              0  &-1 &-5
            \end{mat}
            \grstep{-\rho_2+\rho_3}
            \begin{mat}[r]
              1  &0  &2  \\
              0  &-1 &-5 \\
              0  &0  &0
            \end{mat}
            \grstep{-\rho_2}
            \begin{mat}[r]
              1  &0  &2  \\
              0  &1  &5  \\
              0  &0  &0
            \end{mat}
          \end{equation*}
          Matriks kedua menjadi
          \begin{equation*}
            \grstep{-2\rho_1+\rho_3}
            \begin{mat}[r]
              1  &0  &2  \\
              0  &2  &10 \\
              0  &0  &0
            \end{mat}
            \grstep{(1/2)\rho_2}
            \begin{mat}[r]
              1  &0  &2  \\
              0  &1  &5  \\
              0  &0  &0
            \end{mat}
          \end{equation*}
          Kedua matriks tersebut ekuivalen baris.
        \partsitem Keduanya tidak ekuivalen baris karena ukurannya berbeda.
        \partsitem Matriks pertama menjadi
          \begin{equation*}
            \grstep{\rho_1+\rho_2}
            \begin{mat}[r]
              1  &1  &1  \\
              0  &3  &3
            \end{mat}
            \grstep{(1/3)\rho_2}
            \begin{mat}[r]
              1  &1  &1  \\
              0  &1  &1
            \end{mat}
            \grstep{-\rho_2+\rho_1}
            \begin{mat}[r]
              1  &0  &0  \\
              0  &1  &1
            \end{mat}
          \end{equation*}
          sedangkan matriks kedua menjadi
          \begin{equation*}
            \grstep{\rho_1\leftrightarrow\rho_2}
            \begin{mat}[r]
              2  &2  &5  \\
              0  &3  &-1
            \end{mat}
            \grstep[(1/3)\rho_2]{(1/2)\rho_1}
            \begin{mat}[r]
              1  &1  &5/2 \\
              0  &1  &-1/3
            \end{mat}
            \grstep{-\rho_2+\rho_1}
            \begin{mat}[r]
              1  &0  &17/6 \\
              0  &1  &-1/3
            \end{mat}
          \end{equation*}
          Keduanya tidak ekuivalen baris.
        \partsitem Di sini matriks pertama menjadi
          \begin{equation*}
            \grstep{(1/3)\rho_2}
            \begin{mat}[r]
              1  &1  &1  \\
              0  &0  &1
            \end{mat}
            \grstep{-\rho_2+\rho_1}
            \begin{mat}[r]
              1  &1  &0  \\
              0  &0  &1
            \end{mat}
          \end{equation*}
          sedangkan matriks kedua menjadi
          \begin{equation*}
            \grstep{\rho_1\leftrightarrow\rho_2}
            \begin{mat}[r]
              1  &-1 &1  \\
              0  &1  &2
            \end{mat}
            \grstep{\rho_2+\rho_1}
            \begin{mat}[r]
              1  &0  &3  \\
              0  &1  &2
            \end{mat}
          \end{equation*}
          Keduanya tidak ekuivalen baris.
       \end{exparts}  
     \end{answer}
  \item 
    Manakah di antara matriks-matriks berikut yang saling ekuivalen baris?
    \begin{exparts*}
      \partsitem
         \(\begin{mat}
             1 &3 \\
             2 &4
           \end{mat}\)
      \partsitem
         \(\begin{mat}
             1 &5 \\
             2 &10
           \end{mat}\)
      \partsitem
         \(\begin{mat}
             1 &-1 \\
             3 &0
           \end{mat}\)
      \partsitem
         \(\begin{mat}
             2 &6 \\
             4 &10
           \end{mat}\)
      \partsitem
         \(\begin{mat}
             0  &1 \\
             -1 &0
           \end{mat}\)
      \partsitem
         \(\begin{mat}
             3 &3 \\
             2 &2
           \end{mat}\)
    \end{exparts*}
    \begin{answer}
      Lakukan reduksi Gauss--Jordan pada setiap matriks.
      Dua matriks ekuivalen baris jika dan hanya jika keduanya memiliki bentuk
      eselon tereduksi yang sama.
      Berikut bentuk tereduksi masing-masing matriks.
      % sage: M = matrix(QQ, [[1,3], [2,4]])
      % sage: gauss_jordan(M)
      % [1 3]
      % [2 4]
      %  take -2 times row 1 plus row 2
      % [ 1  3]
      % [ 0 -2]
      %  take -1/2 times row 2
      % [1 3]
      % [0 1]
      %  take -3 times row 2 plus row 1
      % [1 0]
      % [0 1]
      % sage: M = matrix(QQ, [[1,5], [2,10]])
      % sage: gauss_jordan(M)
      % [ 1  5]
      % [ 2 10]
      %  take -2 times row 1 plus row 2
      % [1 5]
      % [0 0]
      % sage: M = matrix(QQ, [[1,-1], [3,0]])
      % sage: gauss_jordan(M)
      % [ 1 -1]
      % [ 3  0]
      %  take -3 times row 1 plus row 2
      % [ 1 -1]
      % [ 0  3]
      %  take 1/3 times row 2
      % [ 1 -1]
      % [ 0  1]
      %  take 1 times row 2 plus row 1
      % [1 0]
      % [0 1]
      % sage: M = matrix(QQ, [[2,6], [4,10]])
      % sage: gauss_jordan(M)
      % [ 2  6]
      % [ 4 10]
      %  take -2 times row 1 plus row 2
      % [ 2  6]
      % [ 0 -2]
      %  take 1/2 times row 1
      %  take -1/2 times row 2
      % [1 3]
      % [0 1]
      %  take -3 times row 2 plus row 1
      % [1 0]
      % [0 1]
      % sage: M = matrix(QQ, [[0,1], [-1,0]])
      % sage: gauss_jordan(M)
      % [ 0  1]
      % [-1  0]
      %  swap row 1 with row 2
      % [-1  0]
      % [ 0  1]
      %  take -1 times row 1
      % [1 0]
      % [0 1]
      % sage: M = matrix(QQ, [[3,3], [2,2]])
      % sage: gauss_jordan(M)
      % [3 3]
      % [2 2]
      %  take -2/3 times row 1 plus row 2
      % [3 3]
      % [0 0]
      %  take 1/3 times row 1
      % [1 1]
      % [0 0]
      \begin{exparts*}
        \partsitem
          \(\begin{mat}
              1  &0  \\
              0  &1
            \end{mat} \)
        \partsitem
          \(\begin{mat}
              1  &5  \\
              0  &0
            \end{mat} \)
        \partsitem
          \(\begin{mat}
             1   &0  \\
             0   &1
            \end{mat} \)
        \partsitem
          \(\begin{mat}
             1   &0  \\
             0   &1
            \end{mat} \)
        \partsitem
          \(\begin{mat}
             1   &0  \\
             0   &1
            \end{mat} \)
        \partsitem
          \(\begin{mat}
             1   &1  \\
             0   &0
            \end{mat} \)
      \end{exparts*}
      Jadi, matriks pada butir (a), (c), (d), dan (e) saling ekuivalen baris.
      Matriks pada butir (b) dan (f) masing-masing berada dalam kelas lain yang
      berbeda.
    \end{answer}
 \item Buatlah tiga matriks lain yang ekuivalen baris dengan matriks yang
   diberikan.
   \begin{exparts*}
     \partsitem 
       \(\begin{mat}
           1 &3 \\
           4 &-1
         \end{mat}\)
     \partsitem
       \(\begin{mat}
           0 &1 &2 \\
           1 &1 &1  \\
           2 &3 &4
         \end{mat}\)
   \end{exparts*}
   \begin{answer}
     Untuk masing-masing matriks, kita cukup melakukan beberapa operasi baris
     pada matriks awal.
     \begin{exparts}
       \partsitem
         Mengalikan baris pertama dengan~$3$ menghasilkan
         \begin{equation*}
           \begin{mat}
             3 &9  \\
             4 &-1
           \end{mat}
         \end{equation*}
         (Pemilihan operasi baris ini tidak memiliki alasan khusus; operasi
         tersebut hanya yang pertama terpikirkan.)
         Dua operasi baris lainnya adalah pertukaran
         $\rho_1\leftrightarrow\rho_2$ dan penambahan $\rho_1+\rho_2$.
         \begin{equation*}
           \begin{mat}
             4 &-1  \\
             1 &3
           \end{mat}
           \quad
           \begin{mat}
             1 &3  \\
             5 &2
           \end{mat}
         \end{equation*}
       \partsitem Menjalankan tiga operasi baris sebarang yang sama menghasilkan
         ketiga matriks berikut.
         \begin{equation*}
          \begin{mat} 
           0 &3 &6 \\
           1 &1 &1  \\
           2 &3 &4  
          \end{mat}
          \quad         
          \begin{mat} 
           1 &1 &1  \\
           0 &1 &2 \\
           2 &3 &4  
          \end{mat}
          \quad         
          \begin{mat} 
           0 &1 &2 \\
           1 &2 &3  \\
           2 &3 &4  
          \end{mat}
          \quad         
         \end{equation*}
     \end{exparts}
   \end{answer}
  \recommended \item Terapkan Metode Gauss pada matriks berikut.
    Nyatakan setiap baris matriks akhir sebagai kombinasi linear baris-baris
    matriks awal.
    \begin{equation*}
      \begin{mat}
        1 &2   &1 \\
        3 &-1  &0 \\
        0 &4   &0
      \end{mat}
    \end{equation*}
    \begin{answer}
      Reduksi Gauss berjalan seperti biasa.
      % sage: load("gauss_method.sage")
      % sage: M = matrix (QQ, [[1,2,1], [3,-1,0], [0,4,0]])
      % sage: gauss_method(M)
      % [ 1  2  1]
      % [ 3 -1  0]
      % [ 0  4  0]
      %  take -3 times row 1 plus row 2
      % [ 1  2  1]
      % [ 0 -7 -3]
      % [ 0  4  0]
      %  take 4/7 times row 2 plus row 3
      % [    1     2     1]
      % [    0    -7    -3]
      % [    0     0 -12/7]
      \begin{equation*}
        \begin{mat}
          1 &2   &1 \\
          3 &-1  &0 \\
          0 &4   &0
        \end{mat}
        \grstep{-3\rho_1+\rho_2}    
        \begin{mat}
          1 &2   &1 \\
          0 &-7  &-3 \\
          0 &4   &0
        \end{mat}
        \grstep{(4/7)\rho_2+\rho_3}    
        \begin{mat}
          1 &2   &1 \\
          0 &-7  &-3 \\
          0 &0   &-12/7
        \end{mat}
      \end{equation*}
      Dengan menamai matriks-matriks tersebut berturut-turut $A$, $D$, dan~$B$,
      kita memperoleh
      \begin{align*}
        \begin{mat}
          \alpha_1 \\
          \alpha_2 \\
          \alpha_3
        \end{mat}
        &\grstep{-3\rho_1+\rho_2}    
        \begin{mat}
          \delta_1=\alpha_1 \\
          \delta_2=-3\alpha_1+\alpha_2 \\
          \delta_3=\alpha_3
        \end{mat}                                   \\
        & \grstep{(4/7)\rho_2+\rho_3}
        \begin{mat}
          \beta_1=\alpha_1  \\
          \beta_2=-3\alpha_1+\alpha_2 \\
          \beta_3=-(12/7)\alpha_1+(4/7)\alpha_2+\alpha_3
        \end{mat}
      \end{align*}
    \end{answer}
  \item 
     Jelaskan matriks-matriks dalam setiap kelas yang wakilnya ditampilkan dalam
     \nearbyexample{ex:RowEqClassTwoTwoMats}.
     \begin{answer}
       Pertama, satu-satunya matriks yang ekuivalen baris dengan matriks yang
       seluruh entrinya \( 0 \) adalah matriks itu sendiri (karena operasi
       baris tidak berpengaruh).

       Kedua, matriks-matriks yang tereduksi menjadi
       \begin{equation*}
         \begin{mat}
           1  &a  \\
           0  &0
         \end{mat}
       \end{equation*}
       berbentuk
       \begin{equation*}
         \begin{mat}
           b  &ba \\
           c  &ca
         \end{mat}
       \end{equation*}
       (dengan \( a,b,c\in\Re \), dan \(b\) serta \(c\) tidak keduanya nol).

       Selanjutnya, matriks-matriks yang tereduksi menjadi
       \begin{equation*}
         \begin{mat}[r]
           0  &1  \\
           0  &0
         \end{mat}
       \end{equation*}
       berbentuk
       \begin{equation*}
         \begin{mat}
           0  &a \\
           0  &b
         \end{mat}
       \end{equation*}
       (dengan \( a,b\in\Re \), dan keduanya tidak sekaligus nol).

       Terakhir, matriks-matriks yang tereduksi menjadi
       \begin{equation*}
         \begin{mat}[r]
           1  &0  \\
           0  &1
         \end{mat}
       \end{equation*}
       adalah matriks-matriks nonsingular.
       Alasannya, sistem linear yang memiliki matriks tersebut sebagai matriks
       koefisien akan memiliki solusi tunggal, dan itulah definisi nonsingular.
       (Cara lain menyatakan hal yang sama adalah bahwa matriks-matriks itu
       tidak termasuk kelas mana pun di atas.)
     \end{answer}
  \item 
    Jelaskan semua matriks dalam kelas ekuivalensi baris masing-masing matriks
    berikut.
    \begin{exparts*}
       \partsitem  \(
           \begin{mat}[r]
             1  &0  \\
             0  &0
           \end{mat}  \)
       \partsitem  \(
           \begin{mat}[r]
             1  &2      \\
             2  &4
           \end{mat} \)
       \partsitem \(
           \begin{mat}[r]
             1  &1      \\
             1  &3
           \end{mat} \)
    \end{exparts*}
    \begin{answer}
      \begin{exparts}
        \partsitem Matriks-matriks tersebut berbentuk
          \begin{equation*}
            \begin{mat}
              a  &0  \\
              b  &0
            \end{mat}
          \end{equation*}
          dengan sedikitnya satu dari \( a,b\in\Re \) bernilai tak nol.
        \partsitem Matriks-matriks tersebut berbentuk
          \begin{equation*}
            \begin{mat}
             a  &2a \\
             b  &2b
            \end{mat}
          \end{equation*}
          dengan sedikitnya satu dari \( a,b\in\Re \) bernilai tak nol.
        \partsitem Matriks yang diberikan nonsingular.
          Jadi, kelas ekuivalensi barisnya terdiri atas semua matriks
          nonsingular $\nbyn{2}$.
          (Sebagai rumus, matriks-matriks itu berbentuk
          \begin{equation*}
            \begin{mat}
              a  &b  \\
              c  &d
            \end{mat}
          \end{equation*}
          untuk \( a,b,c,d\in\Re \), dengan \( ad-bc\neq 0 \).
          Dalam Bab Empat, kita akan melihat bahwa rumus ini menentukan kapan
          matriks \( \nbyn{2} \) bersifat nonsingular.)
      \end{exparts}  
    \end{answer}
  \item 
    Berapa banyak kelas ekuivalensi baris yang ada?
    \begin{answer}
       Tak hingga banyak.
       Sebagai contoh, dalam
       \begin{equation*}
         \begin{mat}
           1  &k  \\
           0  &0
         \end{mat}
       \end{equation*}
       setiap $k\in\Re$ memberikan kelas yang berbeda.
    \end{answer}
  \item 
    Dapatkah suatu kelas ekuivalensi baris memuat matriks dengan ukuran
    berbeda?
    \begin{answer}
      Tidak.
      Operasi baris tidak mengubah ukuran matriks.
    \end{answer}
  \item 
    Seberapa besar kelas-kelas ekuivalensi baris?
    \begin{exparts} 
      \partsitem Tunjukkan bahwa untuk setiap matriks nol, kelasnya berhingga.
      \partsitem Adakah kelas lain yang hanya memuat berhingga banyak anggota?
    \end{exparts}
    \begin{answer}
     \begin{exparts}
      \partsitem Operasi baris pada matriks nol tidak berpengaruh.
        Jadi, setiap matriks tersebut merupakan satu-satunya anggota kelas
        ekuivalensi barisnya.
      \partsitem Tidak.
        Setiap entri tak nol dapat diskalakan ulang.
     \end{exparts}
    \end{answer}
  \recommended \item 
    Berikan dua matriks berbentuk eselon tereduksi yang entri-entri utamanya
    berada dalam kolom-kolom yang sama, tetapi tidak ekuivalen baris.
    \begin{answer}
      Berikut dua contohnya.
      \begin{equation*}
        \begin{mat}[r]
          1  &1  &0  \\
          0  &0  &1
        \end{mat}
        \quad\text{dan}\quad
        \begin{mat}[r]
          1  &0  &0  \\
          0  &0  &1
        \end{mat}
      \end{equation*}  
     \end{answer}
  \recommended \item 
    Tunjukkan bahwa sebarang dua matriks nonsingular \( \nbyn{n} \) ekuivalen
    baris.
    Apakah sebarang dua matriks singular ekuivalen baris?
    \begin{answer}
      Sebarang dua matriks nonsingular \( \nbyn{n} \) memiliki bentuk eselon
      tereduksi yang sama, yaitu matriks yang seluruh entrinya \( 0 \), kecuali
      entri-entri \( 1 \) pada diagonal.
      \begin{equation*}
        \begin{mat}
          1  &0  &       &0  \\
          0  &1  &       &0  \\
             &   &\ddots &   \\
          0  &0  &       &1
        \end{mat}
      \end{equation*}

      Dua matriks singular berukuran sama belum tentu ekuivalen baris.
      Sebagai contoh, kedua matriks singular \( \nbyn{2} \) berikut tidak
      ekuivalen baris.
      \begin{equation*}
        \begin{mat}[r]
          1  &1  \\
          0  &0
        \end{mat}
        \quad\text{dan}\quad
        \begin{mat}[r]
          1  &0  \\
          0  &0
        \end{mat}
      \end{equation*}  
    \end{answer}
  \recommended \item 
    Jelaskan semua kelas ekuivalensi baris yang memuat matriks-matriks berikut.
    \begin{exparts*}
      \partsitem matriks \( \nbyn{2} \)
      \partsitem matriks \( \nbym{2}{3} \)
      \partsitem matriks \( \nbym{3}{2} \)
      \partsitem matriks \( \nbyn{3} \)
    \end{exparts*}
    \begin{answer}
      Karena terdapat tepat satu matriks berbentuk eselon tereduksi dalam
      setiap kelas, kita cukup mendaftarkan matriks-matriks berbentuk eselon
      tereduksi yang mungkin.

      Untuk daftar tersebut, lihat jawaban
      \nearbyexercise{exer:PossRedEchFrms}.
    \end{answer}
  \item  
     \begin{exparts}
          \partsitem Tunjukkan bahwa suatu vektor $\vec{\beta}_0$ merupakan
            kombinasi linear dari anggota-anggota himpunan
            $\set{\vec{\beta}_1,\ldots,\vec{\beta}_n}$ jika dan hanya jika
            terdapat suatu hubungan linear
            $\zero=c_0\vec{\beta}_0+\cdots+c_n\vec{\beta}_n$
            dengan $c_0$ tak nol.
            (\textit{Petunjuk.} Perhatikan kasus
            $\vec{\beta}_0=\zero$.)
         \partsitem Gunakan hasil itu untuk menyederhanakan pembuktian
            \nearbylemma{le:EchFormNoLinCombo}.
       \end{exparts}
       \begin{answer}
          \begin{exparts}
           \partsitem Jika terdapat suatu hubungan linear dengan $c_0$ tak nol,
             kita dapat mengurangkan $c_0\vec{\beta}_0$ dari kedua ruas lalu
             membagi dengan $-c_0$ untuk memperoleh $\vec{\beta}_0$ sebagai
             kombinasi linear dari vektor-vektor lainnya.
             (Catatan: jika tidak ada vektor lain dalam himpunan itu\Dash
             misalnya, jika hubungannya adalah
             $\zero=3\cdot\zero$\Dash pernyataan tersebut tetap benar karena,
             berdasarkan definisi, vektor nol merupakan jumlah kosong, yaitu
             penjumlahan atas himpunan vektor kosong.)

             Sebaliknya, jika $\vec{\beta}_0$ merupakan kombinasi dari
             vektor-vektor lainnya,
             $\vec{\beta}_0=c_1\vec{\beta}_1+\dots+c_n\vec{\beta}_n$
             maka mengurangkan $\vec{\beta}_0$ dari kedua ruas menghasilkan
             hubungan dengan sekurang-kurangnya satu koefisien tak nol, yaitu
             $-1$ di depan $\vec{\beta}_0$.
           \partsitem Baris pertama bukan kombinasi linear dari baris-baris
             lainnya karena alasan yang diberikan dalam pembuktian: pada
             persamaan komponen yang berasal dari kolom berisi entri utama
             baris pertama, satu-satunya entri tak nol adalah entri utama
             baris pertama, sehingga koefisiennya harus nol.
             Jadi, menurut bagian sebelumnya dalam latihan ini, baris pertama
             tidak berada dalam hubungan linear apa pun dengan baris-baris
             lainnya.

             Oleh karena itu, ketika meninjau apakah baris kedua dapat berada
             dalam hubungan linear dengan baris-baris lainnya, kita dapat
             mengabaikan baris pertama. Namun, argumen yang baru saja diterapkan
             pada baris pertama kini berlaku pada baris kedua.
             (Dengan kata lain, kita berargumen menggunakan induksi.)
         \end{exparts}
      \end{answer}
  % \item 
  %   Why, in the proof of \nearbytheorem{th:ReducedEchelonFormIsUnique},
  %   do we bother to restrict to the nonzero rows?
  %   Why not just stick to the relationship that we began with,
  %   $\beta_i=c_{i,1}\delta_1+\dots+c_{i,m}\delta_m$, with $m$ instead of $r$,
  %   and argue using it that the only nonzero coefficient
  %   is \( c_{i,i}  \), which is \( 1 \)?
  %   \begin{answer}
  %      The zero rows could have nonzero coefficients, and
  %      so the statement would not be true.
  %   \end{answer}
  \recommended \item 
   \cite{Trono}
   Tiga pengemudi truk singgah di sebuah kedai tepi jalan.
   Pengemudi pertama membeli empat roti lapis, secangkir kopi, dan sepuluh
   donat seharga \$$8.45$.
   Pengemudi lainnya membeli tiga roti lapis, secangkir kopi, dan tujuh donat
   seharga \$$6.30$.
   Berapa harga yang dibayar pengemudi ketiga untuk satu roti lapis,
   secangkir kopi, dan satu donat?
   \begin{answer}
     Kita mengetahui bahwa $4s+c+10d=8.45$ dan $3s+c+7d=6.30$, dan kita ingin
     mengetahui nilai $s+c+d$.
     Untungnya, $s+c+d$ merupakan kombinasi linear dari $4s+c+10d$ dan
     $3s+c+7d$.
     Dengan menyebut harga yang belum diketahui itu sebagai $p$, kita
     memperoleh reduksi berikut.
     \begin{align*}
       \begin{amat}{3}
         4  &1  &10  &8.45 \\
         3  &1  &7   &6.30 \\
         1  &1  &1   &p
       \end{amat}
       &\grstep[-(1/4)\rho_1+\rho_3]{-(3/4)\rho_1+\rho_2}
       \begin{amat}{3}
         4  &1    &10     &8.45      \\
         0  &1/4  &-1/2   &-0.037\,5 \\
         0  &3/4  &-3/2   &p-2.112\,5
       \end{amat}                                      \\
       &\grstep{-3\rho_2+\rho_3}
       \begin{amat}{3}
         4  &1    &10     &8.45      \\
         0  &1/4  &-1/2   &-0.037\,5 \\
         0  &0    &0      &p-2.00
       \end{amat}
     \end{align*}
     Harga yang dibayar adalah \$$2.00$.
   \end{answer}
  % \item
  %  The fact that Gaussian reduction disallows multiplication of
  %  a row by zero is needed for the proof of uniqueness of reduced echelon form,
  %  or else every matrix would
  %  be row equivalent to a matrix of all zeros.
  %  Where is it used?
  %  \begin{answer}
  %    If multiplication of a row by zero were allowed then
  %    \nearbylemma{le:EquivMatsSameForm}
  %    would not hold.
  %    That is, where
  %    \begin{equation*}
  %      \begin{mat}[r]
  %        1  &3  \\
  %        2  &1
  %      \end{mat}
  %      \grstep{0\rho_2}
  %      \begin{mat}[r]
  %        1  &3  \\
  %        0  &0
  %      \end{mat}
  %    \end{equation*}
  %    all the rows of the second matrix can be expressed as linear combinations
  %    of the rows of the first, but the converse does not hold.
  %    The second row of the first matrix is not a linear combination of the
  %    rows of the second matrix.  
  %  \end{answer}
  \item 
   Lema Kombinasi Linear menyatakan persamaan mana yang dapat diperoleh dari
   reduksi Gauss suatu sistem linear tertentu.
   \begin{enumerate}
     \item Berikan suatu persamaan yang tidak diimplikasikan oleh sistem ini.
       \begin{equation*}
         \begin{linsys}{2}
           3x  &+  &4y  &=  &8 \\
           2x  &+  & y  &=  &3 
         \end{linsys}
       \end{equation*}
     \item Apakah dari setiap sistem tak konsisten dapat diturunkan sembarang
       persamaan?
   \end{enumerate}
   \begin{answer}
     \begin{enumerate}
        \item Jawaban yang mudah adalah berikut ini.
          \begin{equation*}
            0=3
          \end{equation*}
          Untuk jawaban yang tidak terlalu berseloroh, selesaikan sistemnya:
          \begin{equation*}
            \begin{amat}[r]{2}
              3  &4  &8  \\
              2  &1   &3
            \end{amat}
            \grstep{-(2/3)\rho_1+\rho_2}
            \begin{amat}[r]{2}
              3  &4    &8    \\
              0  &-5/3 &-7/3
            \end{amat}
          \end{equation*}
          diperoleh \( y=7/5 \) dan \( x=4/5 \).
          Kini persamaan apa pun yang tidak dipenuhi oleh \( (4/5,7/5) \)
          dapat digunakan, misalnya \( 5x+5y=10 \).
        \item Tidak selalu. Operasi baris hanya membentuk kombinasi linear
          dengan koefisien skalar tetap dari persamaan-persamaan semula.
          Sebagai contoh, sistem tak konsisten yang hanya memuat
          \( 0=5 \) dapat menghasilkan persamaan lain yang ruas kirinya nol,
          tetapi tidak dapat menghasilkan koefisien baru untuk $x$ atau $y$;
          jadi sistem itu tidak dapat menghasilkan \( 3x+2y=4 \).

          Suatu sistem tak konsisten dapat menghasilkan setiap persamaan hanya
          jika baris-baris koefisiennya merentang semua kemungkinan koefisien
          variabel. Ketakkonsistenan sendiri tidak menjamin syarat tersebut.
     \end{enumerate}  
    \end{answer}
  \item 
    \cite{HoffmanKunze}
    Perluas definisi ekuivalensi baris ke sistem linear.
    Menurut definisi Anda, apakah sistem-sistem yang ekuivalen memiliki
    himpunan solusi yang sama?
    \begin{answer}
      Definisikan dua sistem linear sebagai ekuivalen jika matriks
      teraugmentasi keduanya ekuivalen baris.
      Pembuktian bahwa sistem-sistem ekuivalen memiliki himpunan solusi yang
      sama adalah mudah.
    \end{answer}
  \item 
    Dalam matriks berikut,
    \begin{equation*}
      \begin{mat}[r]
        1  &2  &3  \\
        3  &0  &3  \\
        1  &4  &5
      \end{mat}
    \end{equation*}
    jumlah kolom pertama dan kedua sama dengan kolom ketiga.
    \begin{exparts}
      \partsitem Tunjukkan bahwa pernyataan tersebut tetap benar setelah
        operasi baris apa pun.
      \partsitem Buatlah sebuah konjektur.
      \partsitem Buktikan bahwa konjektur tersebut benar.
    \end{exparts}
    \begin{answer}
      \begin{exparts}
        \partsitem Ketiga pertukaran baris yang mungkin mudah diperiksa,
          demikian pula ketiga penskalaan ulang yang mungkin.
          Salah satu dari enam kombinasi baris yang mungkin adalah
          \( k\rho_1+\rho_2 \):
          \begin{equation*}
            \begin{mat}
              1           &2           &3  \\
              k\cdot 1+3  &k\cdot 2+0  &k\cdot 3+3  \\
              1           &4           &5
            \end{mat}
          \end{equation*}
          dan sekali lagi jumlah kolom pertama dan kedua sama dengan kolom
          ketiga. Lima kombinasi lainnya serupa.
        \partsitem Konjektur yang wajar adalah bahwa operasi baris tidak
          mengubah hubungan linear antarkolom.
        \partsitem Pembuktian dengan memeriksa setiap kasus mengikuti sketsa
          yang diberikan pada bagian pertama.
      \end{exparts}  
   \end{answer}
\end{exercises}
